% Complete reviewed Unit 9; editorial master: chapters/ch09/student.md
% Markdown SHA-256: 985f526939b050cefe521fb66300aaaa6cdc188426eda798aaf5f4323d1ca60b
% One-use initial assembly followed by complete content and page review.
\newcommand{\runningtitle}{Unit 9. Complements and reports}
% Layout inherited from the reviewed previous course; set a running title per unit.
% This is a small hand-maintained TeX layout, not a conversion pipeline.
\RequirePackage{fix-cm}
\documentclass[11pt,twoside]{article}
\usepackage{fontspec}
\setmainfont{Cambria}[
  Renderer=HarfBuzz,
  BoldFont=cambriab.ttf,
  ItalicFont=cambriai.ttf,
  BoldItalicFont=cambriaz.ttf
]
\setmonofont{lmmono10-regular.otf}[BoldFont=lmmonolt10-bold.otf]
\usepackage[paperwidth=165mm,paperheight=241mm,inner=21mm,outer=19mm,
  top=22mm,bottom=20mm,headheight=14pt,headsep=7mm,footskip=10mm]{geometry}
\usepackage{xcolor,array,booktabs,longtable,tabularx}
\usepackage{enumitem,fancyhdr}
\usepackage[unicode,hidelinks]{hyperref}
\definecolor{ink}{HTML}{000000}
\definecolor{muted}{HTML}{4C5558}
\setlength{\parindent}{0pt}
\setlength{\parskip}{3.5pt}
\AtBeginDocument{\fontsize{11}{14.3}\selectfont}
\setlength{\tabcolsep}{4pt}
\renewcommand{\arraystretch}{1.10}
\setlist[itemize]{leftmargin=1.3em,label=\textbullet,itemsep=2pt,topsep=3pt,parsep=0pt}
\setlist[enumerate]{leftmargin=1.8em,itemsep=3pt,topsep=4pt,parsep=0pt}
\widowpenalty=10000
\clubpenalty=10000
\emergencystretch=1em
\raggedbottom
\pagestyle{fancy}
\fancyhf{}
\providecommand{\runningtitle}{Old English}
\fancyhead[LE,RO]{\small\itshape \runningtitle}
\fancyhead[LO,RE]{\small Old English}
\fancyfoot[LE,RO]{\small\thepage}
\renewcommand{\headrulewidth}{0.3pt}
\newcommand{\booktitle}[2]{
  {\fontsize{9}{11}\selectfont\color{muted}#1\par}
  \vspace{3pt}
  {\fontsize{20}{24}\selectfont\color{ink}#2\par}
  \vspace{6pt}
}
\newcommand{\source}[1]{\par{\footnotesize\color{muted}#1\par}}
\newcommand{\readingtext}[1]{
  \par\addvspace{5pt}
  \begingroup\leftskip=1.5em
    \fontsize{12}{16}\selectfont #1\par
  \endgroup\addvspace{5pt}
}
\makeatletter
\renewcommand{\section}{\@startsection{section}{1}{\z@}
  {-2.4ex \@plus -0.6ex \@minus -0.2ex}{0.8ex \@plus 0.2ex}
  {\normalfont\fontsize{12}{15}\selectfont\bfseries\color{ink}}}
\renewcommand{\subsection}{\@startsection{subsection}{2}{\z@}
  {-1.8ex \@plus -0.5ex \@minus -0.2ex}{0.5ex \@plus 0.2ex}
  {\normalfont\fontsize{11}{14}\selectfont\bfseries\color{ink}}}
\makeatother

\hyphenpenalty=10000
\exhyphenpenalty=10000
\emergencystretch=3em
\newcommand{\itempara}[2]{\par\addvspace{3pt}\noindent\hangindent=1.8em\hangafter=1\makebox[1.8em][l]{#1.}#2\par}
\newcommand{\reserve}[1]{\par\penalty-100\begingroup\dimen0=\pagegoal\advance\dimen0 by-\pagetotal\ifdim\dimen0<#1\newpage\fi\endgroup}
\begin{document}
\addfontfeatures{Ligatures=TeXOff}
\raggedright
\booktitle{UNIT 9}{Complements and reported content}

% mdparagraph:P001:begin
A clause can state what someone knows, report what someone says, or ask a question within another sentence. An infinitive can also describe an action that someone commands or observes. This unit distinguishes those constructions and adds the vocabulary of occupations used in Cook II.9–12.
% mdparagraph:P001:end

% mdparagraph:P002:begin
The examples and “The fisherman and his companion” are new teaching material. “The Choice of Occupations” is an exact source dialogue with its own help and original notes.
% mdparagraph:P002:end

% mdparagraph:P003:begin
In the tables, 1, 2 and 3 mean first, second and third person. Singular and plural cells are shown even when their spellings coincide. Dual subjects take plural finite forms. Noun and adjective tables give nominative, genitive, dative and accusative. Forms labelled \textbf{reference} are for recognition. Keep the written vowel lengths and the distinction between e and ę.
% mdparagraph:P003:end

\reserve{7\baselineskip}
% mdparagraph:P004:begin
\section*{1. Finite content clauses}
% mdparagraph:P004:end

% teaching: NEW-U09-GRAM-u09-content-that
% teaching: NEW-U09-VOC-thaet-conjunction
% teaching: NEW-U09-VOC-witan
% teaching: NEW-U09-VOC-secgan
% mdparagraph:P005:begin
A \textbf{content clause} states the content of knowing or saying. After the selected uses of \textbf{witan}, “know”, and \textbf{sęcgan}, “say, tell”, it may begin with invariant \textbf{ðæt}, “that”. The clause has its own subject and finite verb.
% mdparagraph:P005:end

% mdparagraph:P006:begin
This conjunction has no case, gender or number. Distinguish it from the inflected \textbf{ðæt} used as a demonstrative or article and from relative \textbf{ðæt}. A relative form has a function inside its clause; the content conjunction introduces the clause without being its subject or object.
% mdparagraph:P006:end

\reserve{8\baselineskip}
% example: NEW-U09-EXAMPLE-01
% mdparagraph:P007:begin
Ic wāt, ðæt se fisc gōd is.
% mdparagraph:P007:end

% mdparagraph:P008:begin
I know that the fish is good.
% mdparagraph:P008:end

% mdparagraph:P009:begin
The speaker presents the fish’s quality as a fact. Se fisc is the subject of is; the whole ðæt clause gives what the speaker knows.
% mdparagraph:P009:end

% teaching: NEW-U09-GRAM-u09-reported-statement
% teaching: NEW-U09-GRAM-u09-report-reference
% mdparagraph:P010:begin
Reported content commonly takes the subjunctive. The narrator presents someone else’s assertion without separately asserting its truth. The subjunctive does not mean that the assertion is false. Indicative content clauses remain possible when the proposition is presented as fact; an exercise that requires a particular mood will state the intended setting.
% mdparagraph:P010:end

\reserve{8\baselineskip}
% example: NEW-U09-EXAMPLE-02
% mdparagraph:P011:begin
Hē sægð, ðæt se fisc gōd sīe.
% mdparagraph:P011:end

% mdparagraph:P012:begin
He says that the fish is good.
% mdparagraph:P012:end

% mdparagraph:P013:begin
Sægð reports the present statement; sīe is present subjunctive. English commonly uses “is” in both this translation and the preceding one.
% mdparagraph:P013:end

\reserve{8\baselineskip}
% example: NEW-U09-EXAMPLE-03
% mdparagraph:P014:begin
Hē sægde, ðæt se fisc gōd wǣre.
% mdparagraph:P014:end

% mdparagraph:P015:begin
He said that the fish was good.
% mdparagraph:P015:end

% mdparagraph:P016:begin
Sægde places the speech in the past. Wǣre is preterite subjunctive: here the fish’s quality belongs to the time of that past speech.
% mdparagraph:P016:end

% mdparagraph:P017:begin
Identify who each pronoun denotes when converting direct speech into a report. The original speaker’s “I” may become “he” if another person reports the words; it need not change if the same person is still speaking. Choose person and tense from the actual setting. The present and past examples above give a useful contrast for simultaneous reported content, not an automatic rule that every tense moves backward after a past verb.
% mdparagraph:P017:end

% mdparagraph:P018:begin
The finite verb often comes later, sometimes finally, in a dependent clause. The examples use that common order. Case endings and context remain decisive when a source uses another arrangement.
% mdparagraph:P018:end

\reserve{7\baselineskip}
% mdparagraph:P019:begin
\section*{2. Dependent questions}
% mdparagraph:P019:end

% teaching: NEW-U09-GRAM-u09-dependent-question
% mdparagraph:P020:begin
A \textbf{dependent question} gives the question whose answer someone knows or requests. It can follow \textbf{witan} or \textbf{sęcgan}. A question word such as \textbf{hwā}, \textbf{hwæt} or \textbf{hwǣm} introduces it.
% mdparagraph:P020:end

% mdparagraph:P021:begin
Choose the question word’s case from its function within the dependent clause: nominative subject, accusative object, dative recipient or other governed complement, and genitive possessor. Its case does not come from the verb of the main clause. A verb such as helpan still requires its dative complement inside a dependent question.
% mdparagraph:P021:end

\reserve{8\baselineskip}
% example: NEW-U09-EXAMPLE-04
% mdparagraph:P022:begin
Ic wāt, hwæt ðū hæbbe.
% mdparagraph:P022:end

% mdparagraph:P023:begin
I know what you have.
% mdparagraph:P023:end

% mdparagraph:P024:begin
Hwæt is the accusative object of hæbbe. Ðū is the subject, so the verb is second-person singular subjunctive even though its spelling matches other singular subjunctive forms.
% mdparagraph:P024:end

% teaching: NEW-U09-GRAM-u09-dependent-whether
% teaching: NEW-U09-GRAM-u09-dative-speaking-recipient
% teaching: NEW-U09-VOC-hwaether-conjunction
% mdparagraph:P025:begin
Invariant \textbf{hwæðer}, “whether”, introduces a dependent yes/\allowbreak{}no question. It differs from the earlier inflected adjective \textbf{hwæðer}, “which of two?”. The conjunction does not agree with a following noun.
% mdparagraph:P025:end

% mdparagraph:P026:begin
With \textbf{sęcgan}, the person addressed is dative. \textbf{Saga mē}, “tell me”, can introduce a question about what happened or whether something is so. This dative addressee differs from the accusative actor in the command construction below.
% mdparagraph:P026:end

\reserve{8\baselineskip}
% example: NEW-U09-EXAMPLE-05
% mdparagraph:P027:begin
Saga mē, hwæðer ðā męnn gearwe sīen.
% mdparagraph:P027:end

% mdparagraph:P028:begin
Tell me whether the men are ready.
% mdparagraph:P028:end

% mdparagraph:P029:begin
Mē names the person to be told. Hwæðer introduces the whole question; ðā męnn is the plural subject of sīen.
% mdparagraph:P029:end

% mdparagraph:P030:begin
Subjunctive is a normal pattern in dependent questions, and is the selected form when the exercises request it. An indicative may be appropriate when the dependent question concerns an asserted fact. The distinction between a dependent question and an ordinary relative clause is one of meaning and construction, not merely the position of a verb.
% mdparagraph:P030:end

\reserve{7\baselineskip}
% mdparagraph:P031:begin
\section*{3. Infinitive complements and object predicates}
% mdparagraph:P031:end

\reserve{7\baselineskip}
% mdparagraph:P032:begin
\subsection*{Commands and their actors}
% mdparagraph:P032:end

% teaching: NEW-U09-GRAM-u09-command-infinitive
% teaching: NEW-U09-VOC-hatan
% mdparagraph:P033:begin
\textbf{Hātan}, “command”, takes an \textbf{accusative actor + plain infinitive}. The actor is the person or thing commanded to perform the infinitive’s action. The infinitive has no person ending and no preceding tō. Its own object retains the case required by its verb.
% mdparagraph:P033:end

\reserve{8\baselineskip}
% example: NEW-U09-EXAMPLE-06
% mdparagraph:P034:begin
Se mann hēt hine ðone fisc beran.
% mdparagraph:P034:end

% mdparagraph:P035:begin
The man ordered him to carry the fish.
% mdparagraph:P035:end

% mdparagraph:P036:begin
Se mann is the subject of hēt. Hine is the commanded actor of beran; ðone fisc is the object carried. Two accusatives need not have the same function.
% mdparagraph:P036:end

\reserve{8\baselineskip}
% example: NEW-U09-EXAMPLE-10
% mdparagraph:P037:begin
Ic wille ðone fisc beran.
% mdparagraph:P037:end

% mdparagraph:P038:begin
I want to carry the fish.
% mdparagraph:P038:end

% mdparagraph:P039:begin
Here ic is both the subject of wille and the understood actor of beran. The earlier same-subject modal pattern does not need a separate accusative actor.
% mdparagraph:P039:end

% mdparagraph:P040:begin
The same finite form of hātan can also mean “call, name”. Compare the complement that follows it. \textbf{Cyning} means “king”; its complete noun pattern is given below.
% mdparagraph:P040:end

\reserve{8\baselineskip}
% example: NEW-U09-EXAMPLE-11
% mdparagraph:P041:begin
Hē hēt ðone mann cyning.
% mdparagraph:P041:end

% mdparagraph:P042:begin
He called the man king.
% mdparagraph:P042:end

% mdparagraph:P043:begin
Cyning is an accusative noun predicate identifying ðone mann. This is the naming construction from Unit 7.
% mdparagraph:P043:end

\reserve{8\baselineskip}
% example: NEW-U09-EXAMPLE-12
% mdparagraph:P044:begin
Hē hēt ðone mann gān.
% mdparagraph:P044:end

% mdparagraph:P045:begin
He ordered the man to go.
% mdparagraph:P045:end

% mdparagraph:P046:begin
Gān is a plain infinitive; ðone mann is its commanded actor. The infinitive marks the command construction.
% mdparagraph:P046:end

\reserve{7\baselineskip}
% mdparagraph:P047:begin
\subsection*{Perceived actions}
% mdparagraph:P047:end

% teaching: NEW-U09-GRAM-u09-perception-infinitive
% teaching: NEW-U09-VOC-geseon
% mdparagraph:P048:begin
\textbf{Gesēon}, “see”, can take an \textbf{accusative actor + plain infinitive} describing the action perceived. The construction says that the observer sees the actor performing that action.
% mdparagraph:P048:end

\reserve{8\baselineskip}
% example: NEW-U09-EXAMPLE-07
% mdparagraph:P049:begin
Ic gesēo ðone mann ðone fisc beran.
% mdparagraph:P049:end

% mdparagraph:P050:begin
I see the man carrying the fish.
% mdparagraph:P050:end

% mdparagraph:P051:begin
Ðone mann is the actor of beran; ðone fisc is its object. English “carrying” translates the Old English infinitive here. It does not make beran a participle.
% mdparagraph:P051:end

\reserve{7\baselineskip}
% mdparagraph:P052:begin
\subsection*{A quality predicated of the object}
% mdparagraph:P052:end

% teaching: NEW-U09-GRAM-u09-object-predicate
% teaching: NEW-U09-VOC-gewyrcean
% mdparagraph:P053:begin
With \textbf{gewyrcean}, “make”, an accusative object may be followed by an adjective predicated of that object. The adjective agrees with the object in case, gender and number. It is strong in this predicate use; an article belonging to the object noun does not make the separate predicate adjective weak.
% mdparagraph:P053:end

\reserve{8\baselineskip}
% example: NEW-U09-EXAMPLE-08
% mdparagraph:P054:begin
God geworhte ðone mann gōdne.
% mdparagraph:P054:end

% mdparagraph:P055:begin
God made the man good.
% mdparagraph:P055:end

% mdparagraph:P056:begin
Gōdne is masculine accusative singular because it describes ðone mann. Its form is not controlled by nominative God.
% mdparagraph:P056:end

\reserve{7\baselineskip}
% mdparagraph:P057:begin
\section*{4. Wishes, preference and related expressions}
% mdparagraph:P057:end

% teaching: NEW-U09-GRAM-u09-unreal-wish
% mdparagraph:P058:begin
The preterite subjunctive can express an unreal wish about the present. The past form need not refer to past time. In such a wish, \textbf{wǣre} means “were” in “would that I were”, and \textbf{hæfde} can mean “had” in “would that I had”. The verb commonly stands first.
% mdparagraph:P058:end

\reserve{8\baselineskip}
% example: NEW-U09-EXAMPLE-09
% mdparagraph:P059:begin
Wǣre ic cyning!
% mdparagraph:P059:end

% mdparagraph:P060:begin
Would that I were king!
% mdparagraph:P060:end

% mdparagraph:P061:begin
The speaker is not king now. This is an unreal present wish, not a statement about a former reign.
% mdparagraph:P061:end

% mdparagraph:P062:begin
Some indicative and subjunctive forms are identical. In particular, first-person singular hæfde can state actual past possession or express an unreal wish. The construction and the stated situation distinguish them; its spelling alone cannot.
% mdparagraph:P062:end

% mdparagraph:P063:begin
The earlier present-subjunctive wishes and exhortations remain available, including first-person plural “let us …” and third-person “let …”. These do not require the English verb “let” to be translated by an additional Old English verb.
% mdparagraph:P063:end

% teaching: NEW-U09-GRAM-u09-volitional-refusal
% teaching: NEW-U09-VOC-willan
% mdparagraph:P064:begin
\textbf{Ne + willan} can express unwillingness or refusal: “does not wish to”, “will not”. Use the familiar positive forms of willan with ne before the finite verb. Here English “will not” can concern someone’s willingness; it need not predict the future. No new contracted negative forms are needed in these exercises.
% mdparagraph:P064:end

% teaching: NEW-U09-GRAM-u09-infinitive-subject
% teaching: NEW-U09-GRAM-u09-mid-accompaniment
% teaching: NEW-U09-VOC-leof-comp
% teaching: NEW-U09-VOC-mid
% mdparagraph:P065:begin
The preference expression \textbf{ūs is lēofre wician …}, “it is preferable to us to sojourn …”, has a dative experiencer, \textbf{ūs}, and the neuter comparative \textbf{lēofre}. The infinitive names what is preferable. \textbf{Wician} means “sojourn, lodge”. English “we had rather” can express this present preference; its word “had” does not require a past verb.
% mdparagraph:P065:end

% mdparagraph:P066:begin
\textbf{Mid + dative} names an accompanying person or cargo. In the preference construction, compare parallel mid phrases on either side of \textbf{ðonne}, “than”: staying with one person is compared with staying with another. Use the dative after each mid. The accompaniment use also covers going aboard with one’s wares.
% mdparagraph:P066:end

% teaching: NEW-U09-GRAM-u09-free-relative-that
% teaching: NEW-U09-VOC-se
% mdparagraph:P067:begin
In the expression “be what you ought to be”, \textbf{ðæt ðe} means “what, that which”, without an expressed noun antecedent. The expression identifies a role or condition and functions as a predicate with “be”. It is distinct from the statement conjunction ðæt.
% mdparagraph:P067:end

% mdparagraph:P068:begin
The familiar \textbf{sculan + plain infinitive} expresses obligation, including “ought to be”. Keep the person and number of sculan appropriate to its subject. This unit uses the predicative ðæt ðe expression in those be/\allowbreak{}ought-to-be clauses.
% mdparagraph:P068:end

% teaching: NEW-U09-GRAM-u09-partitive-pronoun
% teaching: NEW-U09-VOC-hwilc
% teaching: NEW-U09-VOC-sum
% teaching: NEW-U09-VOC-thu
% mdparagraph:P069:begin
A personal genitive can identify the group from which a member is selected: \textbf{hwilc ēower}, “which of you”, and \textbf{sume hiera}, “some of them”. Here \textbf{ēower} is the genitive plural personal pronoun, not an agreeing possessive adjective. The word selecting from the group still has the case required by its own clause.
% mdparagraph:P069:end

% mdparagraph:P070:begin
Cook also gives \textbf{hīe sume}, “some of them”. In this expression hīe and sume are both nominative plural, standing in apposition. Do not parse hīe as genitive. The appositional expression and the genitive construction are two ways to express the selection; Cook’s reference to §151 asks for the former.
% mdparagraph:P070:end

\reserve{7\baselineskip}
% mdparagraph:P071:begin
\section*{5. Nouns: complete patterns and vocabulary}
% mdparagraph:P071:end

% mdparagraph:P072:begin
Each noun is given with its gender. The tables show complete singular and plural patterns. In the vocabulary tables, \textbf{stem} means the base to which the model’s endings are added. Use the headword itself for nominative singular. In the strong masculine and neuter models below, accusative singular is also unchanged; weak masculine and feminine nouns instead follow their model’s -an. The listed genitive singular and nominative plural verify the stem and principal endings.
% mdparagraph:P072:end

% mdparagraph:P073:begin
Each vocabulary group names its full model. Replace that model’s stem by the listed stem, retaining its endings. Do not remove an internal vowel or change a root vowel unless the diagnostic forms show the change. The three new patterns—sunu, rīce and geðeahtend—are given separately in full.
% mdparagraph:P073:end

% mdparagraph:P074:begin
In meanings marked \textbf{singular}, use the noun in that stated material, abstract or collective sense in the singular. Plural counterparts remain reference forms; they do not automatically mean kinds, portions or separate instances. Meanings marked \textbf{plural} use the plural, with singular counterparts for reference. Unmarked nouns are used in both numbers.
% mdparagraph:P074:end

\reserve{7\baselineskip}
% mdparagraph:P075:begin
\subsection*{Strong masculines}
% mdparagraph:P075:end

% mdparagraph:P076:begin
The regular masculine has genitive singular -es and dative -e. Its plural has -as, -a, -um, -as. Cræft keeps æ before its consonant cluster; do not substitute a in its plural.
% mdparagraph:P076:end

% teaching: NEW-U09-NOM-SMITH
\reserve{7\baselineskip}
% mdparagraph:P077:begin
\subsection*{smið}
% mdparagraph:P077:end

% table: NEW-U09-NOM-SMITH
% mdtable:T01:begin
\par\begingroup
\fontsize{10.5}{13.2}\selectfont
\setlength{\LTleft}{0pt}\setlength{\LTright}{\fill}
\noindent
\begin{tabular}{@{}>{\raggedright\arraybackslash}p{39.46mm}>{\raggedright\arraybackslash}p{39.46mm}>{\raggedright\arraybackslash}p{39.46mm}@{}}
\toprule
\textbf{Case} & \textbf{Singular} & \textbf{Plural} \\
\midrule
Nominative & smið & smiðas \\
Genitive & smiðes & smiða \\
Dative & smiðe & smiðum \\
Accusative & smið & smiðas \\
\bottomrule
\end{tabular}
\par\endgroup

% mdtable:T01:end
% teaching: NEW-U09-NOM-SULHSCEAR
% teaching: NEW-U09-VOC-sulhscear
% teaching: NEW-U09-VOC-smith
% teaching: NEW-U09-NOM-BELG
% teaching: NEW-U09-VOC-belg
% teaching: NEW-U09-NOM-IERTHLING
% teaching: NEW-U09-VOC-ierthling
% teaching: NEW-U09-NOM-MAESSEPREOST
% teaching: NEW-U09-VOC-maessepreost
% teaching: NEW-U09-NOM-CRAEFT
% teaching: NEW-U09-VOC-craeft
% teaching: NEW-U09-NOM-THEOWDOM
% teaching: NEW-U09-VOC-theowdom
% teaching: NEW-U09-NOM-ELP
% teaching: NEW-U09-VOC-elp
% teaching: NEW-U09-NOM-CYNING
% teaching: NEW-U09-VOC-cyning
% teaching: NEW-U09-NOM-MUNUC
% teaching: NEW-U09-VOC-munuc
% teaching: NEW-U09-NOM-DRENC
% teaching: NEW-U09-VOC-drenc
% teaching: NEW-U09-NOM-CEORL
% teaching: NEW-U09-VOC-ceorl
% teaching: NEW-U09-NOM-PLIHT
% teaching: NEW-U09-VOC-pliht
% teaching: NEW-U09-NOM-BEORN
% teaching: NEW-U09-VOC-beorn
% mdparagraph:P078:begin
Vocabulary following \textbf{smið}:
% mdparagraph:P078:end

% mdtable:T02:begin
\par\begingroup
\fontsize{10.5}{13.2}\selectfont
\setlength{\LTleft}{0pt}\setlength{\LTright}{\fill}
\noindent
\begin{tabular}{@{}>{\raggedright\arraybackslash}p{21.3mm}>{\raggedright\arraybackslash}p{20mm}>{\raggedright\arraybackslash}p{22.3mm}>{\raggedright\arraybackslash}p{24.6mm}>{\raggedright\arraybackslash}p{24.55mm}@{}}
\toprule
\textbf{Noun} & \textbf{Gender; meaning} & \textbf{Stem} & \textbf{Genitive singular} & \textbf{Nominative plural} \\
\midrule
sulhscear & m.; plowshare & sulhscear- & sulhsceares & sulhscearas \\
smið & m.; blacksmith & smið- & smiðes & smiðas \\
bęlg & m.; bellows (plural) & bęlg- & bęlges & bęlgas \\
ierðling & m.; farmer; plowman & ierðling- & ierðlinges & ierðlingas \\
mæsseprēost & m.; priest & mæsseprēost- & mæsseprēostes & mæsseprēostas \\
cræft & m.; trade; occupation; craft & cræft- & cræftes & cræftas \\
ðēowdōm & m.; service (singular) & ðēowdōm- & ðēowdōmes & ðēowdōmas \\
elp & m.; elephant & elp- & elpes & elpas \\
cyning & m.; king & cyning- & cyninges & cyningas \\
munuc & m.; monk & munuc- & munuces & munucas \\
dręnc & m.; drink (singular) & dręnc- & dręnces & dręncas \\
ceorl & m.; layman & ceorl- & ceorles & ceorlas \\
pliht & m.; peril; risk & pliht- & plihtes & plihtas \\
beorn & m.; warrior & beorn- & beornes & beornas \\
\bottomrule
\end{tabular}
\par\endgroup

% mdtable:T02:end
\reserve{7\baselineskip}
% mdparagraph:P079:begin
\subsection*{Masculines in -e}
% mdparagraph:P079:end

% mdparagraph:P080:begin
The headword’s final -e remains in nominative and accusative singular. Before another ending use the listed stem: fiscer- + -es gives fisceres. Dative singular ends in -e; the plural has the ordinary masculine endings.
% mdparagraph:P080:end

% teaching: NEW-U09-NOM-FISCERE
\reserve{7\baselineskip}
% mdparagraph:P081:begin
\subsection*{fiscere}
% mdparagraph:P081:end

% table: NEW-U09-NOM-FISCERE
% mdtable:T03:begin
\par\begingroup
\fontsize{10.5}{13.2}\selectfont
\setlength{\LTleft}{0pt}\setlength{\LTright}{\fill}
\noindent
\begin{tabular}{@{}>{\raggedright\arraybackslash}p{39.46mm}>{\raggedright\arraybackslash}p{39.46mm}>{\raggedright\arraybackslash}p{39.46mm}@{}}
\toprule
\textbf{Case} & \textbf{Singular} & \textbf{Plural} \\
\midrule
Nominative & fiscere & fisceras \\
Genitive & fisceres & fiscera \\
Dative & fiscere & fiscerum \\
Accusative & fiscere & fisceras \\
\bottomrule
\end{tabular}
\par\endgroup

% mdtable:T03:end
% teaching: NEW-U09-NOM-MANGERE
% teaching: NEW-U09-VOC-mangere
% teaching: NEW-U09-NOM-ELE
% teaching: NEW-U09-VOC-ele
% teaching: NEW-U09-VOC-fiscere
% teaching: NEW-U09-NOM-SEAMERE
% teaching: NEW-U09-VOC-seamere
% teaching: NEW-U09-NOM-EALDORSCIPE
% teaching: NEW-U09-VOC-ealdorscipe
% teaching: NEW-U09-NOM-LYRE
% teaching: NEW-U09-VOC-lyre
% mdparagraph:P082:begin
Vocabulary following \textbf{fiscere}:
% mdparagraph:P082:end

% mdtable:T04:begin
\par\begingroup
\fontsize{10.5}{13.2}\selectfont
\setlength{\LTleft}{0pt}\setlength{\LTright}{\fill}
\noindent
\begin{tabular}{@{}>{\raggedright\arraybackslash}p{19.5mm}>{\raggedright\arraybackslash}p{29.97mm}>{\raggedright\arraybackslash}p{18.5mm}>{\raggedright\arraybackslash}p{22.39mm}>{\raggedright\arraybackslash}p{22.39mm}@{}}
\toprule
\textbf{Noun} & \textbf{Gender; meaning} & \textbf{Stem} & \textbf{Genitive singular} & \textbf{Nominative plural} \\
\midrule
mangere & m.; merchant & manger- & mangeres & mangeras \\
ęle & m.; oil (singular) & ęl- & ęles & ęlas \\
fiscere & m.; fisherman & fiscer- & fisceres & fisceras \\
sēamere & m.; tailor & sēamer- & sēameres & sēameras \\
ealdorscipe & m.; chief place; primacy (singular) & ealdorscip- & ealdorscipes & ealdorscipas \\
lyre & m.; loss (singular) & lyr- & lyres & lyras \\
\bottomrule
\end{tabular}
\par\endgroup

% mdtable:T04:end
\reserve{7\baselineskip}
% mdparagraph:P083:begin
\subsection*{A masculine with a shortened stem}
% mdparagraph:P083:end

% mdparagraph:P084:begin
Angel loses its internal e before endings: use angl-. The headword remains angel in nominative and accusative singular.
% mdparagraph:P084:end

% teaching: NEW-U09-NOM-ANGEL
\reserve{7\baselineskip}
% mdparagraph:P085:begin
\subsection*{angel}
% mdparagraph:P085:end

% table: NEW-U09-NOM-ANGEL
% mdtable:T05:begin
\par\begingroup
\fontsize{10.5}{13.2}\selectfont
\setlength{\LTleft}{0pt}\setlength{\LTright}{\fill}
\noindent
\begin{tabular}{@{}>{\raggedright\arraybackslash}p{39.46mm}>{\raggedright\arraybackslash}p{39.46mm}>{\raggedright\arraybackslash}p{39.46mm}@{}}
\toprule
\textbf{Case} & \textbf{Singular} & \textbf{Plural} \\
\midrule
Nominative & angel & anglas \\
Genitive & angles & angla \\
Dative & angle & anglum \\
Accusative & angel & anglas \\
\bottomrule
\end{tabular}
\par\endgroup

% mdtable:T05:end
% teaching: NEW-U09-VOC-angel
% mdparagraph:P086:begin
Vocabulary following \textbf{angel}:
% mdparagraph:P086:end

% mdtable:T06:begin
\par\begingroup
\fontsize{10.5}{13.2}\selectfont
\setlength{\LTleft}{0pt}\setlength{\LTright}{\fill}
\noindent
\begin{tabular}{@{}>{\raggedright\arraybackslash}p{19.82mm}>{\raggedright\arraybackslash}p{23.18mm}>{\raggedright\arraybackslash}p{19.82mm}>{\raggedright\arraybackslash}p{23.88mm}>{\raggedright\arraybackslash}p{26.06mm}@{}}
\toprule
\textbf{Noun} & \textbf{Gender; meaning} & \textbf{Stem} & \textbf{Genitive singular} & \textbf{Nominative plural} \\
\midrule
angel & m.; fish-hook & angl- & angles & anglas \\
\bottomrule
\end{tabular}
\par\endgroup

% mdtable:T06:end
\reserve{7\baselineskip}
% mdparagraph:P087:begin
\subsection*{Strong neuters with an endingless plural}
% mdparagraph:P087:end

% mdparagraph:P088:begin
These neuters have the usual -es and -e in genitive and dative singular. Nominative and accusative plural have no ending; genitive and dative plural have -a and -um.
% mdparagraph:P088:end

% teaching: NEW-U09-NOM-HUS
\reserve{7\baselineskip}
% mdparagraph:P089:begin
\subsection*{hūs}
% mdparagraph:P089:end

% table: NEW-U09-NOM-HUS
% mdtable:T07:begin
\par\begingroup
\fontsize{10.5}{13.2}\selectfont
\setlength{\LTleft}{0pt}\setlength{\LTright}{\fill}
\noindent
\begin{tabular}{@{}>{\raggedright\arraybackslash}p{39.46mm}>{\raggedright\arraybackslash}p{39.46mm}>{\raggedright\arraybackslash}p{39.46mm}@{}}
\toprule
\textbf{Case} & \textbf{Singular} & \textbf{Plural} \\
\midrule
Nominative & hūs & hūs \\
Genitive & hūses & hūsa \\
Dative & hūse & hūsum \\
Accusative & hūs & hūs \\
\bottomrule
\end{tabular}
\par\endgroup

% mdtable:T07:end
% teaching: NEW-U09-VOC-hus
% teaching: NEW-U09-NOM-WIN
% teaching: NEW-U09-VOC-win
% teaching: NEW-U09-NOM-GOLD
% teaching: NEW-U09-VOC-gold
% teaching: NEW-U09-NOM-HORS
% teaching: NEW-U09-VOC-hors
% teaching: NEW-U09-NOM-GODSPELL
% teaching: NEW-U09-VOC-godspell
% teaching: NEW-U09-NOM-GESTREON
% teaching: NEW-U09-VOC-gestreon
% teaching: NEW-U09-NOM-GEDEORF
% teaching: NEW-U09-VOC-gedeorf
% teaching: NEW-U09-NOM-BAN
% teaching: NEW-U09-VOC-ban
% teaching: NEW-U09-NOM-FOLC
% teaching: NEW-U09-VOC-folc
% teaching: NEW-U09-NOM-GEFLIT
% teaching: NEW-U09-VOC-geflit
% teaching: NEW-U09-NOM-WIF
% teaching: NEW-U09-VOC-wif
% teaching: NEW-U09-NOM-HLAEST
% teaching: NEW-U09-VOC-hlaest
% teaching: NEW-U09-NOM-AER-COPPER
% teaching: NEW-U09-VOC-aer-copper
% mdparagraph:P090:begin
Vocabulary following \textbf{hūs}:
% mdparagraph:P090:end

% mdtable:T08:begin
\par\begingroup
\fontsize{10.5}{13.2}\selectfont
\setlength{\LTleft}{0pt}\setlength{\LTright}{\fill}
\noindent
\begin{tabular}{@{}>{\raggedright\arraybackslash}p{16.11mm}>{\raggedright\arraybackslash}p{38.17mm}>{\raggedright\arraybackslash}p{16.11mm}>{\raggedright\arraybackslash}p{21.18mm}>{\raggedright\arraybackslash}p{21.18mm}@{}}
\toprule
\textbf{Noun} & \textbf{Gender; meaning} & \textbf{Stem} & \textbf{Genitive singular} & \textbf{Nominative plural} \\
\midrule
hūs & n.; house & hūs- & hūses & hūs \\
wīn & n.; wine (singular) & wīn- & wīnes & wīn \\
gold & n.; gold (singular) & gold- & goldes & gold \\
hors & n.; horse & hors- & horses & hors \\
godspell & n.; Gospel (singular) & godspell- & godspelles & godspell \\
gestrēon & n.; gain; profit (singular) & gestrēon- & gestrēones & gestrēon \\
gedeorf & n.; labor; toil (singular) & gedeorf- & gedeorfes & gedeorf \\
bān & n.; bone & bān- & bānes & bān \\
folc & n.; people, as a collective (singular) & folc- & folces & folc \\
geflīt & n.; dispute; strife & geflīt- & geflītes & geflīt \\
wīf & n.; wife & wīf- & wīfes & wīf \\
hlæst & n.; wares; merchandise; cargo (plural) & hlæst- & hlæstes & hlæst \\
ǣr & n.; copper (singular) & ǣr- & ǣres & ǣr \\
\bottomrule
\end{tabular}
\par\endgroup

% mdtable:T08:end
\reserve{7\baselineskip}
% mdparagraph:P091:begin
\subsection*{Short-stem neuters}
% mdparagraph:P091:end

% mdparagraph:P092:begin
The short stem takes -u in nominative and accusative plural. Tin has a short vowel. In this unit the material “tin” is used in the singular; its plural table is a reference model.
% mdparagraph:P092:end

% teaching: NEW-U09-NOM-TIN
\reserve{7\baselineskip}
% mdparagraph:P093:begin
\subsection*{tin}
% mdparagraph:P093:end

% table: NEW-U09-NOM-TIN
% mdtable:T09:begin
\par\begingroup
\fontsize{10.5}{13.2}\selectfont
\setlength{\LTleft}{0pt}\setlength{\LTright}{\fill}
\noindent
\begin{tabular}{@{}>{\raggedright\arraybackslash}p{39.46mm}>{\raggedright\arraybackslash}p{39.46mm}>{\raggedright\arraybackslash}p{39.46mm}@{}}
\toprule
\textbf{Case} & \textbf{Singular} & \textbf{Plural} \\
\midrule
Nominative & tin & tinu \\
Genitive & tines & tina \\
Dative & tine & tinum \\
Accusative & tin & tinu \\
\bottomrule
\end{tabular}
\par\endgroup

% mdtable:T09:end
% teaching: NEW-U09-VOC-tin
% mdparagraph:P094:begin
Vocabulary following \textbf{tin}:
% mdparagraph:P094:end

% mdtable:T10:begin
\par\begingroup
\fontsize{10.5}{13.2}\selectfont
\setlength{\LTleft}{0pt}\setlength{\LTright}{\fill}
\noindent
\begin{tabular}{@{}>{\raggedright\arraybackslash}p{19.18mm}>{\raggedright\arraybackslash}p{26.08mm}>{\raggedright\arraybackslash}p{19.18mm}>{\raggedright\arraybackslash}p{23.11mm}>{\raggedright\arraybackslash}p{25.22mm}@{}}
\toprule
\textbf{Noun} & \textbf{Gender; meaning} & \textbf{Stem} & \textbf{Genitive singular} & \textbf{Nominative plural} \\
\midrule
tin & n.; tin (singular) & tin- & tines & tinu \\
\bottomrule
\end{tabular}
\par\endgroup

% mdtable:T10:end
\reserve{7\baselineskip}
% mdparagraph:P095:begin
\subsection*{A neuter retaining its internal vowel}
% mdparagraph:P095:end

% mdparagraph:P096:begin
Ðyrel has short y and retains its internal e before endings. Its nominative and accusative plural end in -u.
% mdparagraph:P096:end

% teaching: NEW-U09-NOM-THYREL
\reserve{7\baselineskip}
% mdparagraph:P097:begin
\subsection*{ðyrel}
% mdparagraph:P097:end

% table: NEW-U09-NOM-THYREL
% mdtable:T11:begin
\par\begingroup
\fontsize{10.5}{13.2}\selectfont
\setlength{\LTleft}{0pt}\setlength{\LTright}{\fill}
\noindent
\begin{tabular}{@{}>{\raggedright\arraybackslash}p{39.46mm}>{\raggedright\arraybackslash}p{39.46mm}>{\raggedright\arraybackslash}p{39.46mm}@{}}
\toprule
\textbf{Case} & \textbf{Singular} & \textbf{Plural} \\
\midrule
Nominative & ðyrel & ðyrelu \\
Genitive & ðyreles & ðyrela \\
Dative & ðyrele & ðyrelum \\
Accusative & ðyrel & ðyrelu \\
\bottomrule
\end{tabular}
\par\endgroup

% mdtable:T11:end
% teaching: NEW-U09-VOC-thyrel
% mdparagraph:P098:begin
Vocabulary following \textbf{ðyrel}:
% mdparagraph:P098:end

% mdtable:T12:begin
\par\begingroup
\fontsize{10.5}{13.2}\selectfont
\setlength{\LTleft}{0pt}\setlength{\LTright}{\fill}
\noindent
\begin{tabular}{@{}>{\raggedright\arraybackslash}p{20.21mm}>{\raggedright\arraybackslash}p{21.42mm}>{\raggedright\arraybackslash}p{20.21mm}>{\raggedright\arraybackslash}p{24.35mm}>{\raggedright\arraybackslash}p{26.57mm}@{}}
\toprule
\textbf{Noun} & \textbf{Gender; meaning} & \textbf{Stem} & \textbf{Genitive singular} & \textbf{Nominative plural} \\
\midrule
ðyrel & n.; hole & ðyrel- & ðyreles & ðyrelu \\
\bottomrule
\end{tabular}
\par\endgroup

% mdtable:T12:end
\reserve{7\baselineskip}
% mdparagraph:P099:begin
\subsection*{A neuter with a shortened stem}
% mdparagraph:P099:end

% mdparagraph:P100:begin
Fōdor uses fōdr- before endings. Its ordinary sense here is singular “fodder”; the plural counterparts are reference forms.
% mdparagraph:P100:end

% teaching: NEW-U09-NOM-FODOR
\reserve{7\baselineskip}
% mdparagraph:P101:begin
\subsection*{fōdor}
% mdparagraph:P101:end

% table: NEW-U09-NOM-FODOR
% mdtable:T13:begin
\par\begingroup
\fontsize{10.5}{13.2}\selectfont
\setlength{\LTleft}{0pt}\setlength{\LTright}{\fill}
\noindent
\begin{tabular}{@{}>{\raggedright\arraybackslash}p{39.46mm}>{\raggedright\arraybackslash}p{39.46mm}>{\raggedright\arraybackslash}p{39.46mm}@{}}
\toprule
\textbf{Case} & \textbf{Singular} & \textbf{Plural} \\
\midrule
Nominative & fōdor & fōdru \\
Genitive & fōdres & fōdra \\
Dative & fōdre & fōdrum \\
Accusative & fōdor & fōdru \\
\bottomrule
\end{tabular}
\par\endgroup

% mdtable:T13:end
% teaching: NEW-U09-VOC-fodor
% mdparagraph:P102:begin
Vocabulary following \textbf{fōdor}:
% mdparagraph:P102:end

% mdtable:T14:begin
\par\begingroup
\fontsize{10.5}{13.2}\selectfont
\setlength{\LTleft}{0pt}\setlength{\LTright}{\fill}
\noindent
\begin{tabular}{@{}>{\raggedright\arraybackslash}p{18.75mm}>{\raggedright\arraybackslash}p{28.02mm}>{\raggedright\arraybackslash}p{18.75mm}>{\raggedright\arraybackslash}p{22.59mm}>{\raggedright\arraybackslash}p{24.65mm}@{}}
\toprule
\textbf{Noun} & \textbf{Gender; meaning} & \textbf{Stem} & \textbf{Genitive singular} & \textbf{Nominative plural} \\
\midrule
fōdor & n.; fodder (singular) & fōdr- & fōdres & fōdru \\
\bottomrule
\end{tabular}
\par\endgroup

% mdtable:T14:end
\reserve{7\baselineskip}
% mdparagraph:P103:begin
\subsection*{Strong feminines}
% mdparagraph:P103:end

% mdparagraph:P104:begin
The nominative singular is the headword. The other singular cases end in -e; the selected plural has -a, -a, -um, -a. Feminine nominative and accusative plural -e is a documented alternative to selected -a. Both ǣl “awl” and slęcg “hammer” are feminine.
% mdparagraph:P104:end

% teaching: NEW-U09-NOM-NAEDL
\reserve{7\baselineskip}
% mdparagraph:P105:begin
\subsection*{nǣdl}
% mdparagraph:P105:end

% table: NEW-U09-NOM-NAEDL
% mdtable:T15:begin
\par\begingroup
\fontsize{10.5}{13.2}\selectfont
\setlength{\LTleft}{0pt}\setlength{\LTright}{\fill}
\noindent
\begin{tabular}{@{}>{\raggedright\arraybackslash}p{39.46mm}>{\raggedright\arraybackslash}p{39.46mm}>{\raggedright\arraybackslash}p{39.46mm}@{}}
\toprule
\textbf{Case} & \textbf{Singular} & \textbf{Plural} \\
\midrule
Nominative & nǣdl & nǣdla \\
Genitive & nǣdle & nǣdla \\
Dative & nǣdle & nǣdlum \\
Accusative & nǣdle & nǣdla \\
\bottomrule
\end{tabular}
\par\endgroup

% mdtable:T15:end
% teaching: NEW-U09-VOC-naedl
% teaching: NEW-U09-NOM-AEL
% teaching: NEW-U09-VOC-ael
% teaching: NEW-U09-NOM-SLECG
% teaching: NEW-U09-VOC-slecg
% teaching: NEW-U09-NOM-EORTHTILTH
% teaching: NEW-U09-VOC-eorthtilth
% mdparagraph:P106:begin
Vocabulary following \textbf{nǣdl}:
% mdparagraph:P106:end

% mdtable:T16:begin
\par\begingroup
\fontsize{10.5}{13.2}\selectfont
\setlength{\LTleft}{0pt}\setlength{\LTright}{\fill}
\noindent
\begin{tabular}{@{}>{\raggedright\arraybackslash}p{17.76mm}>{\raggedright\arraybackslash}p{31.52mm}>{\raggedright\arraybackslash}p{17.76mm}>{\raggedright\arraybackslash}p{22.37mm}>{\raggedright\arraybackslash}p{23.35mm}@{}}
\toprule
\textbf{Noun} & \textbf{Gender; meaning} & \textbf{Stem} & \textbf{Genitive singular} & \textbf{Nominative plural} \\
\midrule
nǣdl & f.; needle & nǣdl- & nǣdle & nǣdla \\
ǣl & f.; awl & ǣl- & ǣle & ǣla \\
slęcg & f.; hammer & slęcg- & slęcge & slęcga \\
eorðtilð & f.; agriculture (singular) & eorðtilð- & eorðtilðe & eorðtilða \\
\bottomrule
\end{tabular}
\par\endgroup

% mdtable:T16:end
\reserve{7\baselineskip}
% mdparagraph:P107:begin
\subsection*{Feminines in -nes}
% mdparagraph:P107:end

% mdparagraph:P108:begin
Keep the single s of the headword in nominative singular, but use the stem in -ness- before every ending. These nouns take the same case endings as nǣdl.
% mdparagraph:P108:end

% teaching: NEW-U09-NOM-FORLIDENNES
\reserve{7\baselineskip}
% mdparagraph:P109:begin
\subsection*{forlidennes}
% mdparagraph:P109:end

% table: NEW-U09-NOM-FORLIDENNES
% mdtable:T17:begin
\par\begingroup
\fontsize{10.5}{13.2}\selectfont
\setlength{\LTleft}{0pt}\setlength{\LTright}{\fill}
\noindent
\begin{tabular}{@{}>{\raggedright\arraybackslash}p{31.99mm}>{\raggedright\arraybackslash}p{41.59mm}>{\raggedright\arraybackslash}p{44.79mm}@{}}
\toprule
\textbf{Case} & \textbf{Singular} & \textbf{Plural} \\
\midrule
Nominative & forlidennes & forlidennessa \\
Genitive & forlidennesse & forlidennessa \\
Dative & forlidennesse & forlidennessum \\
Accusative & forlidennesse & forlidennessa \\
\bottomrule
\end{tabular}
\par\endgroup

% mdtable:T17:end
% teaching: NEW-U09-VOC-forlidennes
% teaching: NEW-U09-NOM-RIHTWISNES
% teaching: NEW-U09-VOC-rihtwisnes
% mdparagraph:P110:begin
Vocabulary following \textbf{forlidennes}:
% mdparagraph:P110:end

% mdtable:T18:begin
\par\begingroup
\fontsize{10.5}{13.2}\selectfont
\setlength{\LTleft}{0pt}\setlength{\LTright}{\fill}
\noindent
\begin{tabular}{@{}>{\raggedright\arraybackslash}p{19mm}>{\raggedright\arraybackslash}p{25.5mm}>{\raggedright\arraybackslash}p{22mm}>{\raggedright\arraybackslash}p{23.125mm}>{\raggedright\arraybackslash}p{23.125mm}@{}}
\toprule
\textbf{Noun} & \textbf{Gender; meaning} & \textbf{Stem} & \textbf{Genitive singular} & \textbf{Nominative plural} \\
\midrule
forlidennes & f.; shipwreck & forlidenness- & forlidennesse & forlidennessa \\
rihtwīsnes & f.; righteousness (singular) & rihtwīsness- & rihtwīsnesse & rihtwīsnessa \\
\bottomrule
\end{tabular}
\par\endgroup

% mdtable:T18:end
\reserve{7\baselineskip}
% mdparagraph:P111:begin
\subsection*{Weak masculines}
% mdparagraph:P111:end

% mdparagraph:P112:begin
Weak masculine nominative singular ends in -a. The other singular cases and nominative and accusative plural end in -an; genitive plural has -ena, dative plural -um.
% mdparagraph:P112:end

% teaching: NEW-U09-NOM-FYRSPEARCA
\reserve{7\baselineskip}
% mdparagraph:P113:begin
\subsection*{fȳrspearca}
% mdparagraph:P113:end

% table: NEW-U09-NOM-FYRSPEARCA
% mdtable:T19:begin
\par\begingroup
\fontsize{10.5}{13.2}\selectfont
\setlength{\LTleft}{0pt}\setlength{\LTright}{\fill}
\noindent
\begin{tabular}{@{}>{\raggedright\arraybackslash}p{35.87mm}>{\raggedright\arraybackslash}p{39.46mm}>{\raggedright\arraybackslash}p{43.05mm}@{}}
\toprule
\textbf{Case} & \textbf{Singular} & \textbf{Plural} \\
\midrule
Nominative & fȳrspearca & fȳrspearcan \\
Genitive & fȳrspearcan & fȳrspearcena \\
Dative & fȳrspearcan & fȳrspearcum \\
Accusative & fȳrspearcan & fȳrspearcan \\
\bottomrule
\end{tabular}
\par\endgroup

% mdtable:T19:end
% teaching: NEW-U09-VOC-fyrspearca
% teaching: NEW-U09-NOM-TREOWWYRHTA
% teaching: NEW-U09-VOC-treowwyrhta
% teaching: NEW-U09-NOM-SCEOWYRHTA
% teaching: NEW-U09-VOC-sceowyrhta
% teaching: NEW-U09-NOM-CEMPA
% teaching: NEW-U09-VOC-cempa
% teaching: NEW-U09-NOM-WYRHTA
% teaching: NEW-U09-VOC-wyrhta
\par\begin{minipage}{\linewidth}\raggedright
% mdparagraph:P114:begin
Vocabulary following \textbf{fȳrspearca}:
% mdparagraph:P114:end

% mdtable:T20:begin
\par\begingroup
\fontsize{10.5}{13.2}\selectfont
\setlength{\LTleft}{0pt}\setlength{\LTright}{\fill}
\noindent
\begin{tabular}{@{}>{\raggedright\arraybackslash}p{21.5mm}>{\raggedright\arraybackslash}p{21.57mm}>{\raggedright\arraybackslash}p{20.5mm}>{\raggedright\arraybackslash}p{24.59mm}>{\raggedright\arraybackslash}p{24.59mm}@{}}
\toprule
\textbf{Noun} & \textbf{Gender; meaning} & \textbf{Stem} & \textbf{Genitive singular} & \textbf{Nominative plural} \\
\midrule
fȳrspearca & m.; spark & fȳrspearc- & fȳrspearcan & fȳrspearcan \\
trēowwyrhta & m.; carpenter & trēowwyrht- & trēowwyrhtan & trēowwyrhtan \\
scēowyrhta & m.; shoemaker & scēowyrht- & scēowyrhtan & scēowyrhtan \\
cęmpa & m.; soldier & cęmp- & cęmpan & cęmpan \\
wyrhta & m.; workman; craftsman & wyrht- & wyrhtan & wyrhtan \\
\bottomrule
\end{tabular}
\par\endgroup

% mdtable:T20:end
\end{minipage}\par
\reserve{7\baselineskip}
% mdparagraph:P115:begin
\subsection*{Weak feminines}
% mdparagraph:P115:end

% mdparagraph:P116:begin
Weak feminine nominative singular ends in -e. The other seven cells use the same endings as the weak masculine pattern. Smiððe and sīde follow this pattern.
% mdparagraph:P116:end

% teaching: NEW-U09-NOM-SMITHTHE
\reserve{7\baselineskip}
% mdparagraph:P117:begin
\subsection*{smiððe}
% mdparagraph:P117:end

% table: NEW-U09-NOM-SMITHTHE
% mdtable:T21:begin
\par\begingroup
\fontsize{10.5}{13.2}\selectfont
\setlength{\LTleft}{0pt}\setlength{\LTright}{\fill}
\noindent
\begin{tabular}{@{}>{\raggedright\arraybackslash}p{39.46mm}>{\raggedright\arraybackslash}p{39.46mm}>{\raggedright\arraybackslash}p{39.46mm}@{}}
\toprule
\textbf{Case} & \textbf{Singular} & \textbf{Plural} \\
\midrule
Nominative & smiððe & smiððan \\
Genitive & smiððan & smiððena \\
Dative & smiððan & smiððum \\
Accusative & smiððan & smiððan \\
\bottomrule
\end{tabular}
\par\endgroup

% mdtable:T21:end
% teaching: NEW-U09-VOC-smiththe
% teaching: NEW-U09-NOM-SIDE
% teaching: NEW-U09-VOC-side
% mdparagraph:P118:begin
Vocabulary following \textbf{smiððe}:
% mdparagraph:P118:end

% mdtable:T22:begin
\par\begingroup
\fontsize{10.5}{13.2}\selectfont
\setlength{\LTleft}{0pt}\setlength{\LTright}{\fill}
\noindent
\begin{tabular}{@{}>{\raggedright\arraybackslash}p{19.03mm}>{\raggedright\arraybackslash}p{26.74mm}>{\raggedright\arraybackslash}p{19.03mm}>{\raggedright\arraybackslash}p{22.93mm}>{\raggedright\arraybackslash}p{25.02mm}@{}}
\toprule
\textbf{Noun} & \textbf{Gender; meaning} & \textbf{Stem} & \textbf{Genitive singular} & \textbf{Nominative plural} \\
\midrule
smiððe & f.; smithy & smiðð- & smiððan & smiððan \\
sīde & f.; silk (singular) & sīd- & sīdan & sīdan \\
\bottomrule
\end{tabular}
\par\endgroup

% mdtable:T22:end
\reserve{7\baselineskip}
% mdparagraph:P119:begin
\subsection*{A compound of mann}
% mdparagraph:P119:end

% mdparagraph:P120:begin
The final member follows mann: dative singular and nominative/\allowbreak{}accusative plural have ealdormęnn. The other forms retain ealdormann-.
% mdparagraph:P120:end

% teaching: NEW-U09-NOM-EALDORMANN
\reserve{7\baselineskip}
% mdparagraph:P121:begin
\subsection*{ealdormann}
% mdparagraph:P121:end

% table: NEW-U09-NOM-EALDORMANN
% mdtable:T23:begin
\par\begingroup
\fontsize{10.5}{13.2}\selectfont
\setlength{\LTleft}{0pt}\setlength{\LTright}{\fill}
\noindent
\begin{tabular}{@{}>{\raggedright\arraybackslash}p{34.82mm}>{\raggedright\arraybackslash}p{41.78mm}>{\raggedright\arraybackslash}p{41.78mm}@{}}
\toprule
\textbf{Case} & \textbf{Singular} & \textbf{Plural} \\
\midrule
Nominative & ealdormann & ealdormęnn \\
Genitive & ealdormannes & ealdormanna \\
Dative & ealdormęnn & ealdormannum \\
Accusative & ealdormann & ealdormęnn \\
\bottomrule
\end{tabular}
\par\endgroup

% mdtable:T23:end
% teaching: NEW-U09-VOC-ealdormann
% mdparagraph:P122:begin
Vocabulary following \textbf{ealdormann}:
% mdparagraph:P122:end

% mdtable:T24:begin
\par\begingroup
\fontsize{10.5}{13.2}\selectfont
\setlength{\LTleft}{0pt}\setlength{\LTright}{\fill}
\noindent
\begin{tabular}{@{}>{\raggedright\arraybackslash}p{20mm}>{\raggedright\arraybackslash}p{23.1mm}>{\raggedright\arraybackslash}p{21mm}>{\raggedright\arraybackslash}p{25.3mm}>{\raggedright\arraybackslash}p{23.35mm}@{}}
\toprule
\textbf{Noun} & \textbf{Gender; meaning} & \textbf{Stem} & \textbf{Genitive singular} & \textbf{Nominative plural} \\
\midrule
ealdormann & m.; prince; leader; noble & ealdormann- & ealdormannes & ealdormęnn \\
\bottomrule
\end{tabular}
\par\endgroup

% mdtable:T24:end
\reserve{7\baselineskip}
% mdparagraph:P123:begin
\subsection*{Sunu: the u-stem pattern}
% mdparagraph:P123:end

% teaching: NEW-U09-GRAM-u09-sunu-u-stem
% mdparagraph:P124:begin
Sunu has short u. Nominative and accusative singular end in -u; genitive and dative singular in -a. The plural has -a except in the dative, which has -um. Cook also records -u/\allowbreak{}-o variants outside the genitives and dative plural, and later forms sunes/\allowbreak{}sunas following the ordinary masculine pattern. The table gives the selected pattern.
% mdparagraph:P124:end

% teaching: NEW-U09-NOM-SUNU
\reserve{7\baselineskip}
% mdparagraph:P125:begin
\subsection*{sunu}
% mdparagraph:P125:end

% table: NEW-U09-NOM-SUNU
% mdtable:T25:begin
\par\begingroup
\fontsize{10.5}{13.2}\selectfont
\setlength{\LTleft}{0pt}\setlength{\LTright}{\fill}
\noindent
\begin{tabular}{@{}>{\raggedright\arraybackslash}p{39.46mm}>{\raggedright\arraybackslash}p{39.46mm}>{\raggedright\arraybackslash}p{39.46mm}@{}}
\toprule
\textbf{Case} & \textbf{Singular} & \textbf{Plural} \\
\midrule
Nominative & sunu & suna \\
Genitive & suna & suna \\
Dative & suna & sunum \\
Accusative & sunu & suna \\
\bottomrule
\end{tabular}
\par\endgroup

% mdtable:T25:end
% teaching: NEW-U09-VOC-sunu
% mdparagraph:P126:begin
Vocabulary following \textbf{sunu}:
% mdparagraph:P126:end

% mdtable:T26:begin
\par\begingroup
\fontsize{10.5}{13.2}\selectfont
\setlength{\LTleft}{0pt}\setlength{\LTright}{\fill}
\noindent
\begin{tabular}{@{}>{\raggedright\arraybackslash}p{20.21mm}>{\raggedright\arraybackslash}p{21.42mm}>{\raggedright\arraybackslash}p{20.21mm}>{\raggedright\arraybackslash}p{24.35mm}>{\raggedright\arraybackslash}p{26.57mm}@{}}
\toprule
\textbf{Noun} & \textbf{Gender; meaning} & \textbf{Stem} & \textbf{Genitive singular} & \textbf{Nominative plural} \\
\midrule
sunu & m.; son & sun- & suna & suna \\
\bottomrule
\end{tabular}
\par\endgroup

% mdtable:T26:end
\reserve{7\baselineskip}
% mdparagraph:P127:begin
\subsection*{Rīce: a neuter in -e}
% mdparagraph:P127:end

% teaching: NEW-U09-GRAM-u09-rice-neuter-e
% mdparagraph:P128:begin
The final -e remains in the singular where no other ending replaces it. Use rīc- before plural endings: rīcu, rīca, rīcum. Cook also gives forms with inserted i/\allowbreak{}e after c or g; rīciu and rīcia are documented alternatives to selected rīcu and rīca.
% mdparagraph:P128:end

% teaching: NEW-U09-NOM-RICE
\reserve{7\baselineskip}
% mdparagraph:P129:begin
\subsection*{rīce}
% mdparagraph:P129:end

% table: NEW-U09-NOM-RICE
% mdtable:T27:begin
\par\begingroup
\fontsize{10.5}{13.2}\selectfont
\setlength{\LTleft}{0pt}\setlength{\LTright}{\fill}
\noindent
\begin{tabular}{@{}>{\raggedright\arraybackslash}p{39.46mm}>{\raggedright\arraybackslash}p{39.46mm}>{\raggedright\arraybackslash}p{39.46mm}@{}}
\toprule
\textbf{Case} & \textbf{Singular} & \textbf{Plural} \\
\midrule
Nominative & rīce & rīcu \\
Genitive & rīces & rīca \\
Dative & rīce & rīcum \\
Accusative & rīce & rīcu \\
\bottomrule
\end{tabular}
\par\endgroup

% mdtable:T27:end
% teaching: NEW-U09-VOC-rice
% mdparagraph:P130:begin
Vocabulary following \textbf{rīce}:
% mdparagraph:P130:end

% mdtable:T28:begin
\par\begingroup
\fontsize{10.5}{13.2}\selectfont
\setlength{\LTleft}{0pt}\setlength{\LTright}{\fill}
\noindent
\begin{tabular}{@{}>{\raggedright\arraybackslash}p{20.21mm}>{\raggedright\arraybackslash}p{21.42mm}>{\raggedright\arraybackslash}p{20.21mm}>{\raggedright\arraybackslash}p{24.35mm}>{\raggedright\arraybackslash}p{26.57mm}@{}}
\toprule
\textbf{Noun} & \textbf{Gender; meaning} & \textbf{Stem} & \textbf{Genitive singular} & \textbf{Nominative plural} \\
\midrule
rīce & n.; kingdom & rīc- & rīces & rīcu \\
\bottomrule
\end{tabular}
\par\endgroup

% mdtable:T28:end
\reserve{7\baselineskip}
% mdparagraph:P131:begin
\subsection*{Geðeahtend: an agent noun in -end}
% mdparagraph:P131:end

% teaching: NEW-U09-GRAM-u09-agent-end
% mdparagraph:P132:begin
Geðeahtend names a person, a counsellor. Its selected nominative and accusative plural have no ending; genitive plural has -ra and dative plural -um. Plural nominative/\allowbreak{}accusative -e and the rarer -as also occur. This noun pattern is distinct from the adjective inflection of a present participle in -ende.
% mdparagraph:P132:end

% teaching: NEW-U09-NOM-GETHEAHTEND
\reserve{7\baselineskip}
% mdparagraph:P133:begin
\subsection*{geðeahtend}
% mdparagraph:P133:end

% table: NEW-U09-NOM-GETHEAHTEND
% mdtable:T29:begin
\par\begingroup
\fontsize{10.5}{13.2}\selectfont
\setlength{\LTleft}{0pt}\setlength{\LTright}{\fill}
\noindent
\begin{tabular}{@{}>{\raggedright\arraybackslash}p{34.82mm}>{\raggedright\arraybackslash}p{41.78mm}>{\raggedright\arraybackslash}p{41.78mm}@{}}
\toprule
\textbf{Case} & \textbf{Singular} & \textbf{Plural} \\
\midrule
Nominative & geðeahtend & geðeahtend \\
Genitive & geðeahtendes & geðeahtendra \\
Dative & geðeahtende & geðeahtendum \\
Accusative & geðeahtend & geðeahtend \\
\bottomrule
\end{tabular}
\par\endgroup

% mdtable:T29:end
% teaching: NEW-U09-VOC-getheahtend
% mdparagraph:P134:begin
Vocabulary following \textbf{geðeahtend}:
% mdparagraph:P134:end

% mdtable:T30:begin
\par\begingroup
\fontsize{10.5}{13.2}\selectfont
\setlength{\LTleft}{0pt}\setlength{\LTright}{\fill}
\noindent
\begin{tabular}{@{}>{\raggedright\arraybackslash}p{18.59mm}>{\raggedright\arraybackslash}p{22.77mm}>{\raggedright\arraybackslash}p{20.45mm}>{\raggedright\arraybackslash}p{26.49mm}>{\raggedright\arraybackslash}p{24.45mm}@{}}
\toprule
\textbf{Noun} & \textbf{Gender; meaning} & \textbf{Stem} & \textbf{Genitive singular} & \textbf{Nominative plural} \\
\midrule
geðeahtend & m.; counsellor & geðeahtend- & geðeahtendes & geðeahtend \\
\bottomrule
\end{tabular}
\par\endgroup

% mdtable:T30:end
% mdparagraph:P135:begin
Distinguish \textbf{beorn}, “warrior”, from \textbf{cęmpa}, “soldier”. \textbf{Mæsseprēost} is the selected word for “priest”; \textbf{ceorl} here means “layman”. \textbf{Hlæst} is used as plural “wares”. The selected use of \textbf{bęlg} refers to plural bellows/\allowbreak{}bags; choose verb agreement from the Old English number rather than from the shape of the English word “bellows”.
% mdparagraph:P135:end

\reserve{7\baselineskip}
% mdparagraph:P136:begin
\section*{6. Adjectives and agreement}
% mdparagraph:P136:end

% mdparagraph:P137:begin
The complete strong and weak forms follow. Use strong adjectives in ordinary predication or unmarked attribution, weak adjectives after the article/\allowbreak{}demonstrative and in direct address. Comparatives are weak. The tables preserve the selected feminine plural -a; the familiar -e alternative is also valid.
% mdparagraph:P137:end

% mdparagraph:P138:begin
The distinct instrumental is not added to these four-case practice patterns.
% mdparagraph:P138:end

\reserve{7\baselineskip}
% mdparagraph:P139:begin
\subsection*{behēfe — useful}
% mdparagraph:P139:end

% teaching: NEW-U09-VOC-behefe
% mdparagraph:P140:begin
Use behēf- before endings. As an adjective in -e, behēfe retains that e only where another ending does not replace it. The selected strong feminine nominative singular and neuter nominative/\allowbreak{}accusative plural are behēfu.
% mdparagraph:P140:end

% teaching: NEW-U09-ADJ-BEHEFE-STRONG
\reserve{7\baselineskip}
% mdparagraph:P141:begin
\subsection*{behēfe: strong declension, singular}
% mdparagraph:P141:end

% table: NEW-U09-ADJ-BEHEFE-STRONG-SG
% mdtable:T31:begin
\par\begingroup
\fontsize{10.5}{13.2}\selectfont
\setlength{\LTleft}{0pt}\setlength{\LTright}{\fill}
\noindent
\begin{tabular}{@{}>{\raggedright\arraybackslash}p{28.89mm}>{\raggedright\arraybackslash}p{28.89mm}>{\raggedright\arraybackslash}p{28.89mm}>{\raggedright\arraybackslash}p{28.89mm}@{}}
\toprule
\textbf{Case} & \textbf{Masculine} & \textbf{Neuter} & \textbf{Feminine} \\
\midrule
Nominative & behēfe & behēfe & behēfu \\
Genitive & behēfes & behēfes & behēfre \\
Dative & behēfum & behēfum & behēfre \\
Accusative & behēfne & behēfe & behēfe \\
\bottomrule
\end{tabular}
\par\endgroup

% mdtable:T31:end
\reserve{7\baselineskip}
% mdparagraph:P142:begin
\subsection*{behēfe: strong declension, plural}
% mdparagraph:P142:end

% table: NEW-U09-ADJ-BEHEFE-STRONG-PL
% mdtable:T32:begin
\par\begingroup
\fontsize{10.5}{13.2}\selectfont
\setlength{\LTleft}{0pt}\setlength{\LTright}{\fill}
\noindent
\begin{tabular}{@{}>{\raggedright\arraybackslash}p{28.89mm}>{\raggedright\arraybackslash}p{28.89mm}>{\raggedright\arraybackslash}p{28.89mm}>{\raggedright\arraybackslash}p{28.89mm}@{}}
\toprule
\textbf{Case} & \textbf{Masculine} & \textbf{Neuter} & \textbf{Feminine} \\
\midrule
Nominative & behēfe & behēfu & behēfa \\
Genitive & behēfra & behēfra & behēfra \\
Dative & behēfum & behēfum & behēfum \\
Accusative & behēfe & behēfu & behēfa \\
\bottomrule
\end{tabular}
\par\endgroup

% mdtable:T32:end
% teaching: NEW-U09-ADJ-BEHEFE-WEAK
\reserve{7\baselineskip}
% mdparagraph:P143:begin
\subsection*{behēfe: weak declension, singular}
% mdparagraph:P143:end

% table: NEW-U09-ADJ-BEHEFE-WEAK-SG
% mdtable:T33:begin
\par\begingroup
\fontsize{10.5}{13.2}\selectfont
\setlength{\LTleft}{0pt}\setlength{\LTright}{\fill}
\noindent
\begin{tabular}{@{}>{\raggedright\arraybackslash}p{28.89mm}>{\raggedright\arraybackslash}p{28.89mm}>{\raggedright\arraybackslash}p{28.89mm}>{\raggedright\arraybackslash}p{28.89mm}@{}}
\toprule
\textbf{Case} & \textbf{Masculine} & \textbf{Neuter} & \textbf{Feminine} \\
\midrule
Nominative & behēfa & behēfe & behēfe \\
Genitive & behēfan & behēfan & behēfan \\
Dative & behēfan & behēfan & behēfan \\
Accusative & behēfan & behēfe & behēfan \\
\bottomrule
\end{tabular}
\par\endgroup

% mdtable:T33:end
\reserve{7\baselineskip}
% mdparagraph:P144:begin
\subsection*{behēfe: weak declension, plural}
% mdparagraph:P144:end

% table: NEW-U09-ADJ-BEHEFE-WEAK-PL
% mdtable:T34:begin
\par\begingroup
\fontsize{10.5}{13.2}\selectfont
\setlength{\LTleft}{0pt}\setlength{\LTright}{\fill}
\noindent
\begin{tabular}{@{}>{\raggedright\arraybackslash}p{28.89mm}>{\raggedright\arraybackslash}p{28.89mm}>{\raggedright\arraybackslash}p{28.89mm}>{\raggedright\arraybackslash}p{28.89mm}@{}}
\toprule
\textbf{Case} & \textbf{Masculine} & \textbf{Neuter} & \textbf{Feminine} \\
\midrule
Nominative & behēfan & behēfan & behēfan \\
Genitive & behēfra & behēfra & behēfra \\
Dative & behēfum & behēfum & behēfum \\
Accusative & behēfan & behēfan & behēfan \\
\bottomrule
\end{tabular}
\par\endgroup

% mdtable:T34:end
\reserve{7\baselineskip}
% mdparagraph:P145:begin
\subsection*{welig — rich; wealthy}
% mdparagraph:P145:end

% teaching: NEW-U09-VOC-welig
% mdparagraph:P146:begin
Keep the stem welig- in these selected forms, including the medial i. Strong feminine nominative singular and neuter nominative/\allowbreak{}accusative plural have -u: weligu.
% mdparagraph:P146:end

% teaching: NEW-U09-ADJ-WELIG-STRONG
\reserve{7\baselineskip}
% mdparagraph:P147:begin
\subsection*{welig: strong declension, singular}
% mdparagraph:P147:end

% table: NEW-U09-ADJ-WELIG-STRONG-SG
% mdtable:T35:begin
\par\begingroup
\fontsize{10.5}{13.2}\selectfont
\setlength{\LTleft}{0pt}\setlength{\LTright}{\fill}
\noindent
\begin{tabular}{@{}>{\raggedright\arraybackslash}p{28.89mm}>{\raggedright\arraybackslash}p{28.89mm}>{\raggedright\arraybackslash}p{28.89mm}>{\raggedright\arraybackslash}p{28.89mm}@{}}
\toprule
\textbf{Case} & \textbf{Masculine} & \textbf{Neuter} & \textbf{Feminine} \\
\midrule
Nominative & welig & welig & weligu \\
Genitive & weliges & weliges & weligre \\
Dative & weligum & weligum & weligre \\
Accusative & weligne & welig & welige \\
\bottomrule
\end{tabular}
\par\endgroup

% mdtable:T35:end
\reserve{7\baselineskip}
% mdparagraph:P148:begin
\subsection*{welig: strong declension, plural}
% mdparagraph:P148:end

% table: NEW-U09-ADJ-WELIG-STRONG-PL
% mdtable:T36:begin
\par\begingroup
\fontsize{10.5}{13.2}\selectfont
\setlength{\LTleft}{0pt}\setlength{\LTright}{\fill}
\noindent
\begin{tabular}{@{}>{\raggedright\arraybackslash}p{28.89mm}>{\raggedright\arraybackslash}p{28.89mm}>{\raggedright\arraybackslash}p{28.89mm}>{\raggedright\arraybackslash}p{28.89mm}@{}}
\toprule
\textbf{Case} & \textbf{Masculine} & \textbf{Neuter} & \textbf{Feminine} \\
\midrule
Nominative & welige & weligu & weliga \\
Genitive & weligra & weligra & weligra \\
Dative & weligum & weligum & weligum \\
Accusative & welige & weligu & weliga \\
\bottomrule
\end{tabular}
\par\endgroup

% mdtable:T36:end
% teaching: NEW-U09-ADJ-WELIG-WEAK
\reserve{7\baselineskip}
% mdparagraph:P149:begin
\subsection*{welig: weak declension, singular}
% mdparagraph:P149:end

% table: NEW-U09-ADJ-WELIG-WEAK-SG
% mdtable:T37:begin
\par\begingroup
\fontsize{10.5}{13.2}\selectfont
\setlength{\LTleft}{0pt}\setlength{\LTright}{\fill}
\noindent
\begin{tabular}{@{}>{\raggedright\arraybackslash}p{28.89mm}>{\raggedright\arraybackslash}p{28.89mm}>{\raggedright\arraybackslash}p{28.89mm}>{\raggedright\arraybackslash}p{28.89mm}@{}}
\toprule
\textbf{Case} & \textbf{Masculine} & \textbf{Neuter} & \textbf{Feminine} \\
\midrule
Nominative & weliga & welige & welige \\
Genitive & weligan & weligan & weligan \\
Dative & weligan & weligan & weligan \\
Accusative & weligan & welige & weligan \\
\bottomrule
\end{tabular}
\par\endgroup

% mdtable:T37:end
\reserve{7\baselineskip}
% mdparagraph:P150:begin
\subsection*{welig: weak declension, plural}
% mdparagraph:P150:end

% table: NEW-U09-ADJ-WELIG-WEAK-PL
% mdtable:T38:begin
\par\begingroup
\fontsize{10.5}{13.2}\selectfont
\setlength{\LTleft}{0pt}\setlength{\LTright}{\fill}
\noindent
\begin{tabular}{@{}>{\raggedright\arraybackslash}p{28.89mm}>{\raggedright\arraybackslash}p{28.89mm}>{\raggedright\arraybackslash}p{28.89mm}>{\raggedright\arraybackslash}p{28.89mm}@{}}
\toprule
\textbf{Case} & \textbf{Masculine} & \textbf{Neuter} & \textbf{Feminine} \\
\midrule
Nominative & weligan & weligan & weligan \\
Genitive & weligra & weligra & weligra \\
Dative & weligum & weligum & weligum \\
Accusative & weligan & weligan & weligan \\
\bottomrule
\end{tabular}
\par\endgroup

% mdtable:T38:end
\reserve{7\baselineskip}
% mdparagraph:P151:begin
\subsection*{wīs — wise}
% mdparagraph:P151:end

% teaching: NEW-U09-VOC-wis
% mdparagraph:P152:begin
The heavy stem wīs- has no -u in strong feminine nominative singular or neuter nominative/\allowbreak{}accusative plural.
% mdparagraph:P152:end

% teaching: NEW-U09-ADJ-WIS-STRONG
\reserve{7\baselineskip}
% mdparagraph:P153:begin
\subsection*{wīs: strong declension, singular}
% mdparagraph:P153:end

% table: NEW-U09-ADJ-WIS-STRONG-SG
% mdtable:T39:begin
\par\begingroup
\fontsize{10.5}{13.2}\selectfont
\setlength{\LTleft}{0pt}\setlength{\LTright}{\fill}
\noindent
\begin{tabular}{@{}>{\raggedright\arraybackslash}p{28.89mm}>{\raggedright\arraybackslash}p{28.89mm}>{\raggedright\arraybackslash}p{28.89mm}>{\raggedright\arraybackslash}p{28.89mm}@{}}
\toprule
\textbf{Case} & \textbf{Masculine} & \textbf{Neuter} & \textbf{Feminine} \\
\midrule
Nominative & wīs & wīs & wīs \\
Genitive & wīses & wīses & wīsre \\
Dative & wīsum & wīsum & wīsre \\
Accusative & wīsne & wīs & wīse \\
\bottomrule
\end{tabular}
\par\endgroup

% mdtable:T39:end
\reserve{7\baselineskip}
% mdparagraph:P154:begin
\subsection*{wīs: strong declension, plural}
% mdparagraph:P154:end

% table: NEW-U09-ADJ-WIS-STRONG-PL
% mdtable:T40:begin
\par\begingroup
\fontsize{10.5}{13.2}\selectfont
\setlength{\LTleft}{0pt}\setlength{\LTright}{\fill}
\noindent
\begin{tabular}{@{}>{\raggedright\arraybackslash}p{28.89mm}>{\raggedright\arraybackslash}p{28.89mm}>{\raggedright\arraybackslash}p{28.89mm}>{\raggedright\arraybackslash}p{28.89mm}@{}}
\toprule
\textbf{Case} & \textbf{Masculine} & \textbf{Neuter} & \textbf{Feminine} \\
\midrule
Nominative & wīse & wīs & wīsa \\
Genitive & wīsra & wīsra & wīsra \\
Dative & wīsum & wīsum & wīsum \\
Accusative & wīse & wīs & wīsa \\
\bottomrule
\end{tabular}
\par\endgroup

% mdtable:T40:end
% teaching: NEW-U09-ADJ-WIS-WEAK
\reserve{7\baselineskip}
% mdparagraph:P155:begin
\subsection*{wīs: weak declension, singular}
% mdparagraph:P155:end

% table: NEW-U09-ADJ-WIS-WEAK-SG
% mdtable:T41:begin
\par\begingroup
\fontsize{10.5}{13.2}\selectfont
\setlength{\LTleft}{0pt}\setlength{\LTright}{\fill}
\noindent
\begin{tabular}{@{}>{\raggedright\arraybackslash}p{28.89mm}>{\raggedright\arraybackslash}p{28.89mm}>{\raggedright\arraybackslash}p{28.89mm}>{\raggedright\arraybackslash}p{28.89mm}@{}}
\toprule
\textbf{Case} & \textbf{Masculine} & \textbf{Neuter} & \textbf{Feminine} \\
\midrule
Nominative & wīsa & wīse & wīse \\
Genitive & wīsan & wīsan & wīsan \\
Dative & wīsan & wīsan & wīsan \\
Accusative & wīsan & wīse & wīsan \\
\bottomrule
\end{tabular}
\par\endgroup

% mdtable:T41:end
\reserve{7\baselineskip}
% mdparagraph:P156:begin
\subsection*{wīs: weak declension, plural}
% mdparagraph:P156:end

% table: NEW-U09-ADJ-WIS-WEAK-PL
% mdtable:T42:begin
\par\begingroup
\fontsize{10.5}{13.2}\selectfont
\setlength{\LTleft}{0pt}\setlength{\LTright}{\fill}
\noindent
\begin{tabular}{@{}>{\raggedright\arraybackslash}p{28.89mm}>{\raggedright\arraybackslash}p{28.89mm}>{\raggedright\arraybackslash}p{28.89mm}>{\raggedright\arraybackslash}p{28.89mm}@{}}
\toprule
\textbf{Case} & \textbf{Masculine} & \textbf{Neuter} & \textbf{Feminine} \\
\midrule
Nominative & wīsan & wīsan & wīsan \\
Genitive & wīsra & wīsra & wīsra \\
Dative & wīsum & wīsum & wīsum \\
Accusative & wīsan & wīsan & wīsan \\
\bottomrule
\end{tabular}
\par\endgroup

% mdtable:T42:end
\reserve{7\baselineskip}
% mdparagraph:P157:begin
\subsection*{manig — many}
% mdparagraph:P157:end

% teaching: NEW-U09-VOC-manig
% teaching: NEW-U09-GRAM-MANIG
% mdparagraph:P158:begin
Use manig as an ordinary attributive adjective with a plural noun. The singular tables below are reference only; they do not add “many a” or a new substantival use. Keep manig- in the selected forms; neuter nominative/\allowbreak{}accusative plural is manigu.
% mdparagraph:P158:end

% teaching: NEW-U09-ADJ-MANIG-STRONG
\reserve{7\baselineskip}
% mdparagraph:P159:begin
\subsection*{manig: strong declension, singular}
% mdparagraph:P159:end

% table: NEW-U09-ADJ-MANIG-STRONG-SG
% mdtable:T43:begin
\par\begingroup
\fontsize{10.5}{13.2}\selectfont
\setlength{\LTleft}{0pt}\setlength{\LTright}{\fill}
\noindent
\begin{tabular}{@{}>{\raggedright\arraybackslash}p{28.89mm}>{\raggedright\arraybackslash}p{28.89mm}>{\raggedright\arraybackslash}p{28.89mm}>{\raggedright\arraybackslash}p{28.89mm}@{}}
\toprule
\textbf{Case} & \textbf{Masculine} & \textbf{Neuter} & \textbf{Feminine} \\
\midrule
Nominative & manig & manig & manigu \\
Genitive & maniges & maniges & manigre \\
Dative & manigum & manigum & manigre \\
Accusative & manigne & manig & manige \\
\bottomrule
\end{tabular}
\par\endgroup

% mdtable:T43:end
\reserve{7\baselineskip}
% mdparagraph:P160:begin
\subsection*{manig: strong declension, plural}
% mdparagraph:P160:end

% table: NEW-U09-ADJ-MANIG-STRONG-PL
% mdtable:T44:begin
\par\begingroup
\fontsize{10.5}{13.2}\selectfont
\setlength{\LTleft}{0pt}\setlength{\LTright}{\fill}
\noindent
\begin{tabular}{@{}>{\raggedright\arraybackslash}p{28.89mm}>{\raggedright\arraybackslash}p{28.89mm}>{\raggedright\arraybackslash}p{28.89mm}>{\raggedright\arraybackslash}p{28.89mm}@{}}
\toprule
\textbf{Case} & \textbf{Masculine} & \textbf{Neuter} & \textbf{Feminine} \\
\midrule
Nominative & manige & manigu & maniga \\
Genitive & manigra & manigra & manigra \\
Dative & manigum & manigum & manigum \\
Accusative & manige & manigu & maniga \\
\bottomrule
\end{tabular}
\par\endgroup

% mdtable:T44:end
% teaching: NEW-U09-ADJ-MANIG-WEAK
\reserve{7\baselineskip}
% mdparagraph:P161:begin
\subsection*{manig: weak declension, singular}
% mdparagraph:P161:end

% table: NEW-U09-ADJ-MANIG-WEAK-SG
% mdtable:T45:begin
\par\begingroup
\fontsize{10.5}{13.2}\selectfont
\setlength{\LTleft}{0pt}\setlength{\LTright}{\fill}
\noindent
\begin{tabular}{@{}>{\raggedright\arraybackslash}p{28.89mm}>{\raggedright\arraybackslash}p{28.89mm}>{\raggedright\arraybackslash}p{28.89mm}>{\raggedright\arraybackslash}p{28.89mm}@{}}
\toprule
\textbf{Case} & \textbf{Masculine} & \textbf{Neuter} & \textbf{Feminine} \\
\midrule
Nominative & maniga & manige & manige \\
Genitive & manigan & manigan & manigan \\
Dative & manigan & manigan & manigan \\
Accusative & manigan & manige & manigan \\
\bottomrule
\end{tabular}
\par\endgroup

% mdtable:T45:end
\reserve{7\baselineskip}
% mdparagraph:P162:begin
\subsection*{manig: weak declension, plural}
% mdparagraph:P162:end

% table: NEW-U09-ADJ-MANIG-WEAK-PL
% mdtable:T46:begin
\par\begingroup
\fontsize{10.5}{13.2}\selectfont
\setlength{\LTleft}{0pt}\setlength{\LTright}{\fill}
\noindent
\begin{tabular}{@{}>{\raggedright\arraybackslash}p{28.89mm}>{\raggedright\arraybackslash}p{28.89mm}>{\raggedright\arraybackslash}p{28.89mm}>{\raggedright\arraybackslash}p{28.89mm}@{}}
\toprule
\textbf{Case} & \textbf{Masculine} & \textbf{Neuter} & \textbf{Feminine} \\
\midrule
Nominative & manigan & manigan & manigan \\
Genitive & manigra & manigra & manigra \\
Dative & manigum & manigum & manigum \\
Accusative & manigan & manigan & manigan \\
\bottomrule
\end{tabular}
\par\endgroup

% mdtable:T46:end
\reserve{7\baselineskip}
% mdparagraph:P163:begin
\subsection*{wīsra — wiser}
% mdparagraph:P163:end

% teaching: NEW-U09-VOC-wis-comp
% mdparagraph:P164:begin
Add the comparative r to wīs-, without changing the vowel, and use weak endings. Genitive plural wīsra simplifies doubled r after s; do not build wīsrra mechanically. This is the same type of reduction seen in earlier lēofra.
% mdparagraph:P164:end

% teaching: NEW-U09-ADJ-WIS-COMP-WEAK
\reserve{7\baselineskip}
% mdparagraph:P165:begin
\subsection*{wīsra: weak declension, singular}
% mdparagraph:P165:end

% table: NEW-U09-ADJ-WIS-COMP-WEAK-SG
% mdtable:T47:begin
\par\begingroup
\fontsize{10.5}{13.2}\selectfont
\setlength{\LTleft}{0pt}\setlength{\LTright}{\fill}
\noindent
\begin{tabular}{@{}>{\raggedright\arraybackslash}p{28.89mm}>{\raggedright\arraybackslash}p{28.89mm}>{\raggedright\arraybackslash}p{28.89mm}>{\raggedright\arraybackslash}p{28.89mm}@{}}
\toprule
\textbf{Case} & \textbf{Masculine} & \textbf{Neuter} & \textbf{Feminine} \\
\midrule
Nominative & wīsra & wīsre & wīsre \\
Genitive & wīsran & wīsran & wīsran \\
Dative & wīsran & wīsran & wīsran \\
Accusative & wīsran & wīsre & wīsran \\
\bottomrule
\end{tabular}
\par\endgroup

% mdtable:T47:end
\reserve{7\baselineskip}
% mdparagraph:P166:begin
\subsection*{wīsra: weak declension, plural}
% mdparagraph:P166:end

% table: NEW-U09-ADJ-WIS-COMP-WEAK-PL
% mdtable:T48:begin
\par\begingroup
\fontsize{10.5}{13.2}\selectfont
\setlength{\LTleft}{0pt}\setlength{\LTright}{\fill}
\noindent
\begin{tabular}{@{}>{\raggedright\arraybackslash}p{28.89mm}>{\raggedright\arraybackslash}p{28.89mm}>{\raggedright\arraybackslash}p{28.89mm}>{\raggedright\arraybackslash}p{28.89mm}@{}}
\toprule
\textbf{Case} & \textbf{Masculine} & \textbf{Neuter} & \textbf{Feminine} \\
\midrule
Nominative & wīsran & wīsran & wīsran \\
Genitive & wīsra & wīsra & wīsra \\
Dative & wīsrum & wīsrum & wīsrum \\
Accusative & wīsran & wīsran & wīsran \\
\bottomrule
\end{tabular}
\par\endgroup

% mdtable:T48:end
% teaching: NEW-U09-GRAM-USEFUL
% mdparagraph:P167:begin
\textbf{Behēfe}, “useful”, can name the person to whom something is useful in the dative. The adjective itself agrees with the subject; the recipient remains dative regardless of that subject’s gender.
% mdparagraph:P167:end

% teaching: NEW-U09-GRAM-u09-mixed-gender-predicate
% mdparagraph:P168:begin
A predicate adjective referring jointly to nouns of different grammatical genders normally takes the neuter plural. Thus the combination of masculine \textbf{ęle} and neuter \textbf{wīn} calls for the selected neuter plural \textbf{gōd}. The zero ending does not make gōd singular in that use.
% mdparagraph:P168:end

% teaching: NEW-U09-GRAM-ADDRESS
% teaching: NEW-U09-VOC-eala
% teaching: NEW-U09-VOC-EALA
% mdparagraph:P169:begin
In direct address, use the noun’s nominative form and a weak accompanying adjective. The number still matters: address one person with singular forms, more than one with plural forms. \textbf{Ēalā}, “O”, is an invariant interjection that may introduce an address or an exclamation.
% mdparagraph:P169:end

\reserve{7\baselineskip}
% mdparagraph:P170:begin
\section*{7. Verbs: complete finite forms}
% mdparagraph:P170:end

% teaching: NEW-U09-GRAM-u09-new-verb-indicatives
% mdparagraph:P171:begin
Each indicative table gives all three persons in both numbers, in the present and preterite. Plural persons share a form, but still refer to different subjects. Each table includes the verb’s stem changes as well as its endings.
% mdparagraph:P171:end

% mdparagraph:P172:begin
The infinitives heading these tables name the verbs. This unit adds their finite forms; the new infinitive used in a sentence is \textbf{wician} in the preference construction. Earlier productive infinitives such as beran, gān and tēon remain available in the constructions already explained.
% mdparagraph:P172:end

\reserve{7\baselineskip}
% mdparagraph:P173:begin
\subsection*{Weak I}
% mdparagraph:P173:end

\reserve{7\baselineskip}
% mdparagraph:P174:begin
\subsection*{cīepan — sell}
% mdparagraph:P174:end

% teaching: NEW-U09-VOC-ciepan
% mdparagraph:P175:begin
Weak I. The object sold is accusative. The preterite has -pte after p: cīepte.
% mdparagraph:P175:end

% teaching: NEW-U09-PAR-ciepan-IND
\reserve{7\baselineskip}
% mdparagraph:P176:begin
\subsection*{cīepan: complete indicative}
% mdparagraph:P176:end

% table: NEW-U09-PAR-ciepan-IND
% mdtable:T49:begin
\par\begingroup
\fontsize{10.5}{13.2}\selectfont
\setlength{\LTleft}{0pt}\setlength{\LTright}{\fill}
\noindent
\begin{tabular}{@{}>{\raggedright\arraybackslash}p{44.51mm}>{\raggedright\arraybackslash}p{36.93mm}>{\raggedright\arraybackslash}p{36.93mm}@{}}
\toprule
\textbf{Person and number} & \textbf{Present} & \textbf{Preterite} \\
\midrule
1 singular & cīepe & cīepte \\
2 singular & cīepst & cīeptest \\
3 singular & cīepð & cīepte \\
1 plural & cīepað & cīepton \\
2 plural & cīepað & cīepton \\
3 plural & cīepað & cīepton \\
\bottomrule
\end{tabular}
\par\endgroup

% mdtable:T49:end
\reserve{7\baselineskip}
% mdparagraph:P177:begin
\subsection*{āfēdan — feed; support}
% mdparagraph:P177:end

% teaching: NEW-U09-VOC-afedan
% mdparagraph:P178:begin
Weak I. The person supported is accusative. Final d and the third-singular dental ending yield tt in āfētt. The preterite keeps -dde: āfēdde.
% mdparagraph:P178:end

% teaching: NEW-U09-PAR-afedan-IND
\reserve{7\baselineskip}
% mdparagraph:P179:begin
\subsection*{āfēdan: complete indicative}
% mdparagraph:P179:end

% table: NEW-U09-PAR-afedan-IND
% mdtable:T50:begin
\par\begingroup
\fontsize{10.5}{13.2}\selectfont
\setlength{\LTleft}{0pt}\setlength{\LTright}{\fill}
\noindent
\begin{tabular}{@{}>{\raggedright\arraybackslash}p{44.51mm}>{\raggedright\arraybackslash}p{36.93mm}>{\raggedright\arraybackslash}p{36.93mm}@{}}
\toprule
\textbf{Person and number} & \textbf{Present} & \textbf{Preterite} \\
\midrule
1 singular & āfēde & āfēdde \\
2 singular & āfētst & āfēddest \\
3 singular & āfētt & āfēdde \\
1 plural & āfēdað & āfēddon \\
2 plural & āfēdað & āfēddon \\
3 plural & āfēdað & āfēddon \\
\bottomrule
\end{tabular}
\par\endgroup

% mdtable:T50:end
\reserve{7\baselineskip}
% mdparagraph:P180:begin
\subsection*{bycgan — buy}
% mdparagraph:P180:end

% teaching: NEW-U09-VOC-bycgan
% mdparagraph:P181:begin
Weak I with an exceptional preterite. The goods bought are accusative. Keep bycg- in bycge/\allowbreak{}bycgað, but single g in bygest/\allowbreak{}bygeð. The preterite has bohte and bohton; do not form a regular preterite from bycg-.
% mdparagraph:P181:end

% teaching: NEW-U09-PAR-bycgan-IND
\reserve{7\baselineskip}
% mdparagraph:P182:begin
\subsection*{bycgan: complete indicative}
% mdparagraph:P182:end

% table: NEW-U09-PAR-bycgan-IND
% mdtable:T51:begin
\par\begingroup
\fontsize{10.5}{13.2}\selectfont
\setlength{\LTleft}{0pt}\setlength{\LTright}{\fill}
\noindent
\begin{tabular}{@{}>{\raggedright\arraybackslash}p{44.51mm}>{\raggedright\arraybackslash}p{36.93mm}>{\raggedright\arraybackslash}p{36.93mm}@{}}
\toprule
\textbf{Person and number} & \textbf{Present} & \textbf{Preterite} \\
\midrule
1 singular & bycge & bohte \\
2 singular & bygest & bohtest \\
3 singular & bygeð & bohte \\
1 plural & bycgað & bohton \\
2 plural & bycgað & bohton \\
3 plural & bycgað & bohton \\
\bottomrule
\end{tabular}
\par\endgroup

% mdtable:T51:end
\reserve{7\baselineskip}
% mdparagraph:P183:begin
\subsection*{sēcan — seek}
% mdparagraph:P183:end

% teaching: NEW-U09-VOC-secan
% mdparagraph:P184:begin
Weak I with preterite sōhte/\allowbreak{}sōhton. The thing sought is accusative. The present and preterite stems also supply the complete subjunctive below.
% mdparagraph:P184:end

% teaching: NEW-U09-PAR-secan-IND
\reserve{7\baselineskip}
% mdparagraph:P185:begin
\subsection*{sēcan: complete indicative}
% mdparagraph:P185:end

% table: NEW-U09-PAR-secan-IND
% mdtable:T52:begin
\par\begingroup
\fontsize{10.5}{13.2}\selectfont
\setlength{\LTleft}{0pt}\setlength{\LTright}{\fill}
\noindent
\begin{tabular}{@{}>{\raggedright\arraybackslash}p{44.51mm}>{\raggedright\arraybackslash}p{36.93mm}>{\raggedright\arraybackslash}p{36.93mm}@{}}
\toprule
\textbf{Person and number} & \textbf{Present} & \textbf{Preterite} \\
\midrule
1 singular & sēce & sōhte \\
2 singular & sēcst & sōhtest \\
3 singular & sēcð & sōhte \\
1 plural & sēcað & sōhton \\
2 plural & sēcað & sōhton \\
3 plural & sēcað & sōhton \\
\bottomrule
\end{tabular}
\par\endgroup

% mdtable:T52:end
\reserve{7\baselineskip}
% mdparagraph:P186:begin
\subsection*{gestillan — settle; quell}
% mdparagraph:P186:end

% teaching: NEW-U09-VOC-gestillan
% mdparagraph:P187:begin
Weak I. The disturbance or dispute quelled is accusative. The doubled l is reduced before the preterite dental ending: gestilde.
% mdparagraph:P187:end

% teaching: NEW-U09-PAR-gestillan-IND
\reserve{7\baselineskip}
% mdparagraph:P188:begin
\subsection*{gestillan: complete indicative}
% mdparagraph:P188:end

% table: NEW-U09-PAR-gestillan-IND
% mdtable:T53:begin
\par\begingroup
\fontsize{10.5}{13.2}\selectfont
\setlength{\LTleft}{0pt}\setlength{\LTright}{\fill}
\noindent
\begin{tabular}{@{}>{\raggedright\arraybackslash}p{44.51mm}>{\raggedright\arraybackslash}p{36.93mm}>{\raggedright\arraybackslash}p{36.93mm}@{}}
\toprule
\textbf{Person and number} & \textbf{Present} & \textbf{Preterite} \\
\midrule
1 singular & gestille & gestilde \\
2 singular & gestilst & gestildest \\
3 singular & gestilð & gestilde \\
1 plural & gestillað & gestildon \\
2 plural & gestillað & gestildon \\
3 plural & gestillað & gestildon \\
\bottomrule
\end{tabular}
\par\endgroup

% mdtable:T53:end
\reserve{7\baselineskip}
% mdparagraph:P189:begin
\subsection*{fręmman — benefit}
% mdparagraph:P189:end

% teaching: NEW-U09-VOC-fremman
% teaching: NEW-U09-GRAM-u09-benefit-dative
% mdparagraph:P190:begin
Weak I. Keep doubled mm in fręmme/\allowbreak{}fręmmað; the second and third singular have fręmest/\allowbreak{}fręmeð, and the preterite has fręmede. The person benefited is dative.
% mdparagraph:P190:end

% teaching: NEW-U09-PAR-fremman-IND
\reserve{7\baselineskip}
% mdparagraph:P191:begin
\subsection*{fręmman: complete indicative}
% mdparagraph:P191:end

% table: NEW-U09-PAR-fremman-IND
% mdtable:T54:begin
\par\begingroup
\fontsize{10.5}{13.2}\selectfont
\setlength{\LTleft}{0pt}\setlength{\LTright}{\fill}
\noindent
\begin{tabular}{@{}>{\raggedright\arraybackslash}p{44.51mm}>{\raggedright\arraybackslash}p{36.93mm}>{\raggedright\arraybackslash}p{36.93mm}@{}}
\toprule
\textbf{Person and number} & \textbf{Present} & \textbf{Preterite} \\
\midrule
1 singular & fręmme & fręmede \\
2 singular & fręmest & fręmedest \\
3 singular & fręmeð & fręmede \\
1 plural & fręmmað & fręmedon \\
2 plural & fręmmað & fręmedon \\
3 plural & fręmmað & fręmedon \\
\bottomrule
\end{tabular}
\par\endgroup

% mdtable:T54:end
\reserve{7\baselineskip}
% mdparagraph:P192:begin
\subsection*{ālīehtan — illuminate; give light to}
% mdparagraph:P192:end

% teaching: NEW-U09-VOC-aliehtan
% mdparagraph:P193:begin
Weak I. The thing illuminated is accusative. The selected second and third singular forms are uncontracted ālīehtest/\allowbreak{}ālīehteð. The preterite is ālīehte: the dental ending joins the stem’s ht without adding another t. Present first singular and preterite first/\allowbreak{}third singular therefore coincide; present and preterite second singular also coincide.
% mdparagraph:P193:end

% teaching: NEW-U09-PAR-aliehtan-IND
\reserve{7\baselineskip}
% mdparagraph:P194:begin
\subsection*{ālīehtan: complete indicative}
% mdparagraph:P194:end

% table: NEW-U09-PAR-aliehtan-IND
% mdtable:T55:begin
\par\begingroup
\fontsize{10.5}{13.2}\selectfont
\setlength{\LTleft}{0pt}\setlength{\LTright}{\fill}
\noindent
\begin{tabular}{@{}>{\raggedright\arraybackslash}p{44.51mm}>{\raggedright\arraybackslash}p{36.93mm}>{\raggedright\arraybackslash}p{36.93mm}@{}}
\toprule
\textbf{Person and number} & \textbf{Present} & \textbf{Preterite} \\
\midrule
1 singular & ālīehte & ālīehte \\
2 singular & ālīehtest & ālīehtest \\
3 singular & ālīehteð & ālīehte \\
1 plural & ālīehtað & ālīehton \\
2 plural & ālīehtað & ālīehton \\
3 plural & ālīehtað & ālīehton \\
\bottomrule
\end{tabular}
\par\endgroup

% mdtable:T55:end
\reserve{7\baselineskip}
% mdparagraph:P195:begin
\subsection*{rǣdan — read}
% mdparagraph:P195:end

% teaching: NEW-U09-VOC-raedan
% mdparagraph:P196:begin
The selected preterite is weak: rǣdde/\allowbreak{}rǣddon. Present third singular rǣtt shows the joining of the stem dental and ending. Rǣdan can take an accusative text or stand without an expressed object, “read”. Strong preterites rēd/\allowbreak{}rēdon are documented alternatives, for recognition rather than an additional system to recall here.
% mdparagraph:P196:end

% teaching: NEW-U09-PAR-raedan-IND
\reserve{7\baselineskip}
% mdparagraph:P197:begin
\subsection*{rǣdan: complete indicative}
% mdparagraph:P197:end

% table: NEW-U09-PAR-raedan-IND
% mdtable:T56:begin
\par\begingroup
\fontsize{10.5}{13.2}\selectfont
\setlength{\LTleft}{0pt}\setlength{\LTright}{\fill}
\noindent
\begin{tabular}{@{}>{\raggedright\arraybackslash}p{44.51mm}>{\raggedright\arraybackslash}p{36.93mm}>{\raggedright\arraybackslash}p{36.93mm}@{}}
\toprule
\textbf{Person and number} & \textbf{Present} & \textbf{Preterite} \\
\midrule
1 singular & rǣde & rǣdde \\
2 singular & rǣtst & rǣddest \\
3 singular & rǣtt & rǣdde \\
1 plural & rǣdað & rǣddon \\
2 plural & rǣdað & rǣddon \\
3 plural & rǣdað & rǣddon \\
\bottomrule
\end{tabular}
\par\endgroup

% mdtable:T56:end
\reserve{7\baselineskip}
% mdparagraph:P198:begin
\subsection*{Weak II}
% mdparagraph:P198:end

\reserve{7\baselineskip}
% mdparagraph:P199:begin
\subsection*{wician — sojourn; lodge}
% mdparagraph:P199:end

% teaching: NEW-U09-VOC-wician
% mdparagraph:P200:begin
Weak II. The present has -ie, -ast, -að and plural -iað; the preterite has -ode/\allowbreak{}-odest/\allowbreak{}-odon. Use mid + dative for the person with whom one lodges. The plain infinitive wician is also used in the preference expression.
% mdparagraph:P200:end

% teaching: NEW-U09-PAR-wician-IND
\reserve{7\baselineskip}
% mdparagraph:P201:begin
\subsection*{wician: complete indicative}
% mdparagraph:P201:end

% table: NEW-U09-PAR-wician-IND
% mdtable:T57:begin
\par\begingroup
\fontsize{10.5}{13.2}\selectfont
\setlength{\LTleft}{0pt}\setlength{\LTright}{\fill}
\noindent
\begin{tabular}{@{}>{\raggedright\arraybackslash}p{44.51mm}>{\raggedright\arraybackslash}p{36.93mm}>{\raggedright\arraybackslash}p{36.93mm}@{}}
\toprule
\textbf{Person and number} & \textbf{Present} & \textbf{Preterite} \\
\midrule
1 singular & wicie & wicode \\
2 singular & wicast & wicodest \\
3 singular & wicað & wicode \\
1 plural & wiciað & wicodon \\
2 plural & wiciað & wicodon \\
3 plural & wiciað & wicodon \\
\bottomrule
\end{tabular}
\par\endgroup

% mdtable:T57:end
\reserve{7\baselineskip}
% mdparagraph:P202:begin
\subsection*{ðolian — suffer}
% mdparagraph:P202:end

% teaching: NEW-U09-VOC-tholian
% mdparagraph:P203:begin
Weak II, with the same endings as wician. The hardship or loss suffered is accusative.
% mdparagraph:P203:end

% teaching: NEW-U09-PAR-tholian-IND
\reserve{7\baselineskip}
% mdparagraph:P204:begin
\subsection*{ðolian: complete indicative}
% mdparagraph:P204:end

% table: NEW-U09-PAR-tholian-IND
% mdtable:T58:begin
\par\begingroup
\fontsize{10.5}{13.2}\selectfont
\setlength{\LTleft}{0pt}\setlength{\LTright}{\fill}
\noindent
\begin{tabular}{@{}>{\raggedright\arraybackslash}p{44.51mm}>{\raggedright\arraybackslash}p{36.93mm}>{\raggedright\arraybackslash}p{36.93mm}@{}}
\toprule
\textbf{Person and number} & \textbf{Present} & \textbf{Preterite} \\
\midrule
1 singular & ðolie & ðolode \\
2 singular & ðolast & ðolodest \\
3 singular & ðolað & ðolode \\
1 plural & ðoliað & ðolodon \\
2 plural & ðoliað & ðolodon \\
3 plural & ðoliað & ðolodon \\
\bottomrule
\end{tabular}
\par\endgroup

% mdtable:T58:end
\reserve{7\baselineskip}
% mdparagraph:P205:begin
\subsection*{gegaderian — gather}
% mdparagraph:P205:end

% teaching: NEW-U09-VOC-gegaderian
% mdparagraph:P206:begin
Weak II. The thing or things gathered are accusative. Keep the prefix ge- throughout the selected forms.
% mdparagraph:P206:end

% teaching: NEW-U09-PAR-gegaderian-IND
\reserve{7\baselineskip}
% mdparagraph:P207:begin
\subsection*{gegaderian: complete indicative}
% mdparagraph:P207:end

% table: NEW-U09-PAR-gegaderian-IND
% mdtable:T59:begin
\par\begingroup
\fontsize{10.5}{13.2}\selectfont
\setlength{\LTleft}{0pt}\setlength{\LTright}{\fill}
\noindent
\begin{tabular}{@{}>{\raggedright\arraybackslash}p{41.89mm}>{\raggedright\arraybackslash}p{34.77mm}>{\raggedright\arraybackslash}p{41.72mm}@{}}
\toprule
\textbf{Person and number} & \textbf{Present} & \textbf{Preterite} \\
\midrule
1 singular & gegaderie & gegaderode \\
2 singular & gegaderast & gegaderodest \\
3 singular & gegaderað & gegaderode \\
1 plural & gegaderiað & gegaderodon \\
2 plural & gegaderiað & gegaderodon \\
3 plural & gegaderiað & gegaderodon \\
\bottomrule
\end{tabular}
\par\endgroup

% mdtable:T59:end
\reserve{7\baselineskip}
% mdparagraph:P208:begin
\subsection*{Strong verbs}
% mdparagraph:P208:end

\reserve{7\baselineskip}
% mdparagraph:P209:begin
\subsection*{rōwan — row}
% mdparagraph:P209:end

% teaching: NEW-U09-VOC-rowan
% mdparagraph:P210:begin
Strong. The second and third singular present change ō to ē: rēwst/\allowbreak{}rēwð. The preterite has rēow-. The verb can stand without an object; ofer introduces an expressed route across something.
% mdparagraph:P210:end

% teaching: NEW-U09-PAR-rowan-IND
\reserve{7\baselineskip}
% mdparagraph:P211:begin
\subsection*{rōwan: complete indicative}
% mdparagraph:P211:end

% table: NEW-U09-PAR-rowan-IND
% mdtable:T60:begin
\par\begingroup
\fontsize{10.5}{13.2}\selectfont
\setlength{\LTleft}{0pt}\setlength{\LTright}{\fill}
\noindent
\begin{tabular}{@{}>{\raggedright\arraybackslash}p{44.51mm}>{\raggedright\arraybackslash}p{36.93mm}>{\raggedright\arraybackslash}p{36.93mm}@{}}
\toprule
\textbf{Person and number} & \textbf{Present} & \textbf{Preterite} \\
\midrule
1 singular & rōwe & rēow \\
2 singular & rēwst & rēowe \\
3 singular & rēwð & rēow \\
1 plural & rōwað & rēowon \\
2 plural & rōwað & rēowon \\
3 plural & rōwað & rēowon \\
\bottomrule
\end{tabular}
\par\endgroup

% mdtable:T60:end
\reserve{7\baselineskip}
% mdparagraph:P212:begin
\subsection*{blāwan — blow}
% mdparagraph:P212:end

% teaching: NEW-U09-VOC-blawan
% mdparagraph:P213:begin
Strong. The second and third singular present have ǣ; the preterite has blēow-. In the bellows sentence the verb stands without an expressed object.
% mdparagraph:P213:end

% teaching: NEW-U09-PAR-blawan-IND
\reserve{7\baselineskip}
% mdparagraph:P214:begin
\subsection*{blāwan: complete indicative}
% mdparagraph:P214:end

% table: NEW-U09-PAR-blawan-IND
% mdtable:T61:begin
\par\begingroup
\fontsize{10.5}{13.2}\selectfont
\setlength{\LTleft}{0pt}\setlength{\LTright}{\fill}
\noindent
\begin{tabular}{@{}>{\raggedright\arraybackslash}p{44.51mm}>{\raggedright\arraybackslash}p{36.93mm}>{\raggedright\arraybackslash}p{36.93mm}@{}}
\toprule
\textbf{Person and number} & \textbf{Present} & \textbf{Preterite} \\
\midrule
1 singular & blāwe & blēow \\
2 singular & blǣwst & blēowe \\
3 singular & blǣwð & blēow \\
1 plural & blāwað & blēowon \\
2 plural & blāwað & blēowon \\
3 plural & blāwað & blēowon \\
\bottomrule
\end{tabular}
\par\endgroup

% mdtable:T61:end
\reserve{7\baselineskip}
% mdparagraph:P215:begin
\subsection*{forlǣtan — abandon; leave}
% mdparagraph:P215:end

% teaching: NEW-U09-VOC-forlaetan
% mdparagraph:P216:begin
Strong. The person or thing abandoned is accusative. The preterite stem is forlēt-. Present third singular forlǣtt joins the stem t and dental ending.
% mdparagraph:P216:end

% teaching: NEW-U09-PAR-forlaetan-IND
\reserve{7\baselineskip}
% mdparagraph:P217:begin
\subsection*{forlǣtan: complete indicative}
% mdparagraph:P217:end

% table: NEW-U09-PAR-forlaetan-IND
% mdtable:T62:begin
\par\begingroup
\fontsize{10.5}{13.2}\selectfont
\setlength{\LTleft}{0pt}\setlength{\LTright}{\fill}
\noindent
\begin{tabular}{@{}>{\raggedright\arraybackslash}p{44.51mm}>{\raggedright\arraybackslash}p{36.93mm}>{\raggedright\arraybackslash}p{36.93mm}@{}}
\toprule
\textbf{Person and number} & \textbf{Present} & \textbf{Preterite} \\
\midrule
1 singular & forlǣte & forlēt \\
2 singular & forlǣtst & forlēte \\
3 singular & forlǣtt & forlēt \\
1 plural & forlǣtað & forlēton \\
2 plural & forlǣtað & forlēton \\
3 plural & forlǣtað & forlēton \\
\bottomrule
\end{tabular}
\par\endgroup

% mdtable:T62:end
\reserve{7\baselineskip}
% mdparagraph:P218:begin
\subsection*{āstīgan — go aboard; board}
% mdparagraph:P218:end

% teaching: NEW-U09-VOC-astigan
% teaching: NEW-U09-GRAM-BOARD
% mdparagraph:P219:begin
Strong I. Present g becomes h before st/\allowbreak{}ð: āstīhst/\allowbreak{}āstīhð. The selected preterite singular is āstāg, second singular āstige, plural āstigon. Final āstāh is a documented alternative to āstāg. The ship boarded is a direct accusative object.
% mdparagraph:P219:end

% teaching: NEW-U09-PAR-astigan-IND
\reserve{7\baselineskip}
% mdparagraph:P220:begin
\subsection*{āstīgan: complete indicative}
% mdparagraph:P220:end

% table: NEW-U09-PAR-astigan-IND
% mdtable:T63:begin
\par\begingroup
\fontsize{10.5}{13.2}\selectfont
\setlength{\LTleft}{0pt}\setlength{\LTright}{\fill}
\noindent
\begin{tabular}{@{}>{\raggedright\arraybackslash}p{44.51mm}>{\raggedright\arraybackslash}p{36.93mm}>{\raggedright\arraybackslash}p{36.93mm}@{}}
\toprule
\textbf{Person and number} & \textbf{Present} & \textbf{Preterite} \\
\midrule
1 singular & āstīge & āstāg \\
2 singular & āstīhst & āstige \\
3 singular & āstīhð & āstāg \\
1 plural & āstīgað & āstigon \\
2 plural & āstīgað & āstigon \\
3 plural & āstīgað & āstigon \\
\bottomrule
\end{tabular}
\par\endgroup

% mdtable:T63:end
\reserve{7\baselineskip}
% mdparagraph:P221:begin
\subsection*{bēatan — beat}
% mdparagraph:P221:end

% teaching: NEW-U09-VOC-beatan
% mdparagraph:P222:begin
Strong. The person or thing beaten is accusative. Present second and third singular change ēa to īe; bīett joins the stem dental and ending. The preterite has bēot-.
% mdparagraph:P222:end

% teaching: NEW-U09-PAR-beatan-IND
\reserve{7\baselineskip}
% mdparagraph:P223:begin
\subsection*{bēatan: complete indicative}
% mdparagraph:P223:end

% table: NEW-U09-PAR-beatan-IND
% mdtable:T64:begin
\par\begingroup
\fontsize{10.5}{13.2}\selectfont
\setlength{\LTleft}{0pt}\setlength{\LTright}{\fill}
\noindent
\begin{tabular}{@{}>{\raggedright\arraybackslash}p{44.51mm}>{\raggedright\arraybackslash}p{36.93mm}>{\raggedright\arraybackslash}p{36.93mm}@{}}
\toprule
\textbf{Person and number} & \textbf{Present} & \textbf{Preterite} \\
\midrule
1 singular & bēate & bēot \\
2 singular & bīetst & bēote \\
3 singular & bīett & bēot \\
1 plural & bēatað & bēoton \\
2 plural & bēatað & bēoton \\
3 plural & bēatað & bēoton \\
\bottomrule
\end{tabular}
\par\endgroup

% mdtable:T64:end
\reserve{7\baselineskip}
% mdparagraph:P224:begin
\subsection*{The prefixed gān pattern}
% mdparagraph:P224:end

\reserve{7\baselineskip}
% mdparagraph:P225:begin
\subsection*{begān — ply; practise}
% mdparagraph:P225:end

% teaching: NEW-U09-VOC-began
% mdparagraph:P226:begin
Begān follows the earlier gān pattern with a prefix. The occupation practised is accusative. Keep begǣst/\allowbreak{}begǣð in the present and beēode/\allowbreak{}beēodest/\allowbreak{}beēodon in the preterite; the prefix does not turn the verb into a weak verb.
% mdparagraph:P226:end

% teaching: NEW-U09-PAR-began-IND
\reserve{7\baselineskip}
% mdparagraph:P227:begin
\subsection*{begān: complete indicative}
% mdparagraph:P227:end

% table: NEW-U09-PAR-began-IND
% mdtable:T65:begin
\par\begingroup
\fontsize{10.5}{13.2}\selectfont
\setlength{\LTleft}{0pt}\setlength{\LTright}{\fill}
\noindent
\begin{tabular}{@{}>{\raggedright\arraybackslash}p{44.51mm}>{\raggedright\arraybackslash}p{36.93mm}>{\raggedright\arraybackslash}p{36.93mm}@{}}
\toprule
\textbf{Person and number} & \textbf{Present} & \textbf{Preterite} \\
\midrule
1 singular & begā & beēode \\
2 singular & begǣst & beēodest \\
3 singular & begǣð & beēode \\
1 plural & begāð & beēodon \\
2 plural & begāð & beēodon \\
3 plural & begāð & beēodon \\
\bottomrule
\end{tabular}
\par\endgroup

% mdtable:T65:end
\reserve{7\baselineskip}
% mdparagraph:P228:begin
\subsection*{Commands}
% mdparagraph:P228:end

% teaching: NEW-U09-GRAM-u09-new-commands
% mdparagraph:P229:begin
The singular command addresses one person; the plural command addresses more than one. The tables give both forms of these six verbs. Their plural commands have the same spelling as the present indicative plurals, but their use in an instruction is different. The singular commands rōw, byge, sēc, begā, gestil and rǣd are shown with their plural counterparts.
% mdparagraph:P229:end

% teaching: NEW-U09-PAR-rowan-IMP
\reserve{7\baselineskip}
% mdparagraph:P230:begin
\subsection*{rōwan: both commands}
% mdparagraph:P230:end

% table: NEW-U09-PAR-rowan-IMP
% mdtable:T66:begin
\par\begingroup
\fontsize{10.5}{13.2}\selectfont
\setlength{\LTleft}{0pt}\setlength{\LTright}{\fill}
\noindent
\begin{tabular}{@{}>{\raggedright\arraybackslash}p{66.23mm}>{\raggedright\arraybackslash}p{54.96mm}@{}}
\toprule
\textbf{Person and number} & \textbf{Command} \\
\midrule
Second singular & rōw \\
Second plural & rōwað \\
\bottomrule
\end{tabular}
\par\endgroup

% mdtable:T66:end
% teaching: NEW-U09-PAR-bycgan-IMP
\reserve{7\baselineskip}
% mdparagraph:P231:begin
\subsection*{bycgan: both commands}
% mdparagraph:P231:end

% table: NEW-U09-PAR-bycgan-IMP
% mdtable:T67:begin
\par\begingroup
\fontsize{10.5}{13.2}\selectfont
\setlength{\LTleft}{0pt}\setlength{\LTright}{\fill}
\noindent
\begin{tabular}{@{}>{\raggedright\arraybackslash}p{66.23mm}>{\raggedright\arraybackslash}p{54.96mm}@{}}
\toprule
\textbf{Person and number} & \textbf{Command} \\
\midrule
Second singular & byge \\
Second plural & bycgað \\
\bottomrule
\end{tabular}
\par\endgroup

% mdtable:T67:end
% teaching: NEW-U09-PAR-secan-IMP
\reserve{7\baselineskip}
% mdparagraph:P232:begin
\subsection*{sēcan: both commands}
% mdparagraph:P232:end

% table: NEW-U09-PAR-secan-IMP
% mdtable:T68:begin
\par\begingroup
\fontsize{10.5}{13.2}\selectfont
\setlength{\LTleft}{0pt}\setlength{\LTright}{\fill}
\noindent
\begin{tabular}{@{}>{\raggedright\arraybackslash}p{66.23mm}>{\raggedright\arraybackslash}p{54.96mm}@{}}
\toprule
\textbf{Person and number} & \textbf{Command} \\
\midrule
Second singular & sēc \\
Second plural & sēcað \\
\bottomrule
\end{tabular}
\par\endgroup

% mdtable:T68:end
% teaching: NEW-U09-PAR-began-IMP
\reserve{7\baselineskip}
% mdparagraph:P233:begin
\subsection*{begān: both commands}
% mdparagraph:P233:end

% table: NEW-U09-PAR-began-IMP
% mdtable:T69:begin
\par\begingroup
\fontsize{10.5}{13.2}\selectfont
\setlength{\LTleft}{0pt}\setlength{\LTright}{\fill}
\noindent
\begin{tabular}{@{}>{\raggedright\arraybackslash}p{66.23mm}>{\raggedright\arraybackslash}p{54.96mm}@{}}
\toprule
\textbf{Person and number} & \textbf{Command} \\
\midrule
Second singular & begā \\
Second plural & begāð \\
\bottomrule
\end{tabular}
\par\endgroup

% mdtable:T69:end
% teaching: NEW-U09-PAR-gestillan-IMP
\reserve{7\baselineskip}
% mdparagraph:P234:begin
\subsection*{gestillan: both commands}
% mdparagraph:P234:end

% table: NEW-U09-PAR-gestillan-IMP
% mdtable:T70:begin
\par\begingroup
\fontsize{10.5}{13.2}\selectfont
\setlength{\LTleft}{0pt}\setlength{\LTright}{\fill}
\noindent
\begin{tabular}{@{}>{\raggedright\arraybackslash}p{66.23mm}>{\raggedright\arraybackslash}p{54.96mm}@{}}
\toprule
\textbf{Person and number} & \textbf{Command} \\
\midrule
Second singular & gestil \\
Second plural & gestillað \\
\bottomrule
\end{tabular}
\par\endgroup

% mdtable:T70:end
% teaching: NEW-U09-PAR-raedan-IMP
\reserve{7\baselineskip}
% mdparagraph:P235:begin
\subsection*{rǣdan: both commands}
% mdparagraph:P235:end

% table: NEW-U09-PAR-raedan-IMP
% mdtable:T71:begin
\par\begingroup
\fontsize{10.5}{13.2}\selectfont
\setlength{\LTleft}{0pt}\setlength{\LTright}{\fill}
\noindent
\begin{tabular}{@{}>{\raggedright\arraybackslash}p{66.23mm}>{\raggedright\arraybackslash}p{54.96mm}@{}}
\toprule
\textbf{Person and number} & \textbf{Command} \\
\midrule
Second singular & rǣd \\
Second plural & rǣdað \\
\bottomrule
\end{tabular}
\par\endgroup

% mdtable:T71:end
\reserve{7\baselineskip}
% mdparagraph:P236:begin
\subsection*{Sēcan: subjunctive}
% mdparagraph:P236:end

% teaching: NEW-U09-GRAM-u09-secan-subjunctive
% mdparagraph:P237:begin
The present subjunctive uses sēc-, the preterite sōht-. Each has singular -e and plural -en for all three persons. Use the present first-person plural in a “let us seek” exhortation. These forms can also be used in the subjunctive constructions explained above.
% mdparagraph:P237:end

% teaching: NEW-U09-PAR-secan-SUBJ
\reserve{7\baselineskip}
% mdparagraph:P238:begin
\subsection*{sēcan: complete subjunctive}
% mdparagraph:P238:end

% table: NEW-U09-PAR-secan-SUBJ
% mdtable:T72:begin
\par\begingroup
\fontsize{10.5}{13.2}\selectfont
\setlength{\LTleft}{0pt}\setlength{\LTright}{\fill}
\noindent
\begin{tabular}{@{}>{\raggedright\arraybackslash}p{44.51mm}>{\raggedright\arraybackslash}p{36.93mm}>{\raggedright\arraybackslash}p{36.93mm}@{}}
\toprule
\textbf{Person and number} & \textbf{Present} & \textbf{Preterite} \\
\midrule
1 singular & sēce & sōhte \\
2 singular & sēce & sōhte \\
3 singular & sēce & sōhte \\
1 plural & sēcen & sōhten \\
2 plural & sēcen & sōhten \\
3 plural & sēcen & sōhten \\
\bottomrule
\end{tabular}
\par\endgroup

% mdtable:T72:end
\reserve{7\baselineskip}
% mdparagraph:P239:begin
\section*{8. Further uses and small words}
% mdparagraph:P239:end

% teaching: NEW-U09-GRAM-u09-absolute-striking
% teaching: NEW-U09-VOC-slean
% mdparagraph:P240:begin
\textbf{Slēan}, “strike, smite”, can stand without naming an object when it means “make a stroke”. Its earlier use with an accusative object remains available; the absence of an object here does not require a new conjugation.
% mdparagraph:P240:end

% teaching: NEW-U09-GRAM-u09-bring-forth
% teaching: NEW-U09-VOC-teon
% teaching: NEW-U09-VOC-forth
% mdparagraph:P241:begin
Use \textbf{tēon … forð}, “bring forth, produce”, with an accusative thing produced. \textbf{Forð} is an invariant adverb. The familiar plain infinitive tēon can follow commanding hātan; the present subjunctive can express a wish or exhortation. Keep the actor of the action distinct from what it produces.
% mdparagraph:P241:end

% teaching: NEW-U09-GRAM-u09-temporal-on
% teaching: NEW-U09-VOC-on
% mdparagraph:P242:begin
\textbf{On + accusative} places an event on a numbered day. The article and ordinal agree with the accusative day noun. This time expression differs from ordinary \textbf{on + dative} locating something in a place.
% mdparagraph:P242:end

% teaching: NEW-U09-GRAM-u09-causal-for-tham
% teaching: NEW-U09-VOC-for-tham
% teaching: NEW-U09-VOC-hwilum
% teaching: NEW-U09-VOC-HWILUM
% mdparagraph:P243:begin
\textbf{For-ðām}, “because”, introduces a factual reason; use the indicative in the reason clause. \textbf{Hwīlum}, “sometimes”, is an invariant adverb of frequency. Its ending does not require you to supply a noun case construction.
% mdparagraph:P243:end

% teaching: NEW-U09-GRAM-SERVICE
% mdparagraph:P244:begin
\textbf{Godes ðēowdōm} means “service of God”: God is the one served. A genitive relationship need not be simple ownership. This expression uses the new service sense of ðēowdōm.
% mdparagraph:P244:end

\reserve{7\baselineskip}
% mdparagraph:P245:begin
\section*{9. Practice in forms and constructions}
% mdparagraph:P245:end

% mdparagraph:P246:begin
All exercises labelled with a letter are new teaching material. Preserve vowel length and the selected forms in recall tasks. In translation, preserve the stated meaning and requested construction. When a task specifies a mood, a grammatically possible sentence using another mood may still fail that particular target.
% mdparagraph:P246:end

\reserve{7\baselineskip}
% mdparagraph:P247:begin
\subsection*{Forms and case use}
% mdparagraph:P247:end

\reserve{5\baselineskip}
\par\begin{minipage}{\linewidth}\raggedright
% mdparagraph:P248:begin
\textbf{F01.}
% mdparagraph:P248:end

% item: NEW-U09-F01
% mdparagraph:P249:begin
Give the nominative, genitive, dative and accusative singular and plural of sunu. Distinguish the dative singular from the nominative and accusative singular.
% mdparagraph:P249:end

\end{minipage}\par\addvspace{4pt}

\par\begin{minipage}{\linewidth}\raggedright
% mdparagraph:P250:begin
\textbf{F02.}
% mdparagraph:P250:end

% item: NEW-U09-F02
% mdparagraph:P251:begin
Give the complete present and preterite indicative of bycgan. Name all three persons in each number, including cells with identical spellings.
% mdparagraph:P251:end

\end{minipage}\par\addvspace{4pt}

\par\begin{minipage}{\linewidth}\raggedright
% mdparagraph:P252:begin
\textbf{F03.}
% mdparagraph:P252:end

% item: NEW-U09-F03
% mdparagraph:P253:begin
Give the singular and plural commands of rōwan and bycgan. State which forms address one person and which address more than one.
% mdparagraph:P253:end

\end{minipage}\par\addvspace{4pt}

\par\begin{minipage}{\linewidth}\raggedright
% mdparagraph:P254:begin
\textbf{F04.}
% mdparagraph:P254:end

% item: NEW-U09-F04
% mdparagraph:P255:begin
Give the nominative, genitive, dative and accusative singular and plural of rīce and geðeahtend. Identify the characteristic plural endings of each pattern.
% mdparagraph:P255:end

\end{minipage}\par\addvspace{4pt}

\par\begin{minipage}{\linewidth}\raggedright
% mdparagraph:P256:begin
\textbf{F05.}
% mdparagraph:P256:end

% item: NEW-U09-F05
% mdparagraph:P257:begin
Translate the nominative noun phrase ‘the kingdoms of the counsellors’, placing the genitive dependent first and omitting a separate article for the head noun. Then translate ‘I give the fish (sing.) to the counsellors.’ Explain the different cases of ‘counsellors’.
% mdparagraph:P257:end

\end{minipage}\par\addvspace{4pt}

\reserve{7\baselineskip}
% mdparagraph:P258:begin
\subsection*{Content and report}
% mdparagraph:P258:end

\reserve{5\baselineskip}
\par\begin{minipage}{\linewidth}\raggedright
% mdparagraph:P259:begin
\textbf{A01.}
% mdparagraph:P259:end

% item: NEW-U09-A01
% mdparagraph:P260:begin
Classify every ðæt in these sentences as demonstrative/\allowbreak{}article, relative, or content conjunction. For each inflected occurrence, give its case and explain whether it agrees with a noun or functions within the relative clause. (a) Ðæt scip is gōd. (b) Ic wāt, ðæt ðæt scip gōd is. (c) Ðæt scip, ðæt ðū hæfst, is gōd.
% mdparagraph:P260:end

\end{minipage}\par\addvspace{4pt}

\par\begin{minipage}{\linewidth}\raggedright
% mdparagraph:P261:begin
\textbf{A02.}
% mdparagraph:P261:end

% item: NEW-U09-A02
% mdparagraph:P262:begin
Translate: I know that the ships are empty. Present their emptiness as an asserted fact and use the indicative in the content clause.
% mdparagraph:P262:end

\end{minipage}\par\addvspace{4pt}

\par\begin{minipage}{\linewidth}\raggedright
% mdparagraph:P263:begin
\textbf{A03.}
% mdparagraph:P263:end

% item: NEW-U09-A03
% mdparagraph:P264:begin
Se fiscere cwið: “Ic hæbbe twēgen fiscas.” Report his words beginning Se fiscere sægð and using ðæt with the present subjunctive. The fisherman is the possessor in both versions.
% mdparagraph:P264:end

\end{minipage}\par\addvspace{4pt}

\par\begin{minipage}{\linewidth}\raggedright
% mdparagraph:P265:begin
\textbf{A04.}
% mdparagraph:P265:end

% item: NEW-U09-A04
% mdparagraph:P266:begin
Se mann cwæð: “Ic hæbbe twēgen fiscas.” Report the same state at the time of his past speech, beginning Se mann sægde. Use ðæt and the preterite subjunctive; the same man possessed the fish.
% mdparagraph:P266:end

\end{minipage}\par\addvspace{4pt}

\par\begin{minipage}{\linewidth}\raggedright
% mdparagraph:P267:begin
\textbf{A05.}
% mdparagraph:P267:end

% item: NEW-U09-A05
% mdparagraph:P268:begin
A task requests a subjunctive report, but a learner writes Hē sægð, ðæt hē gōd is. Is the sentence necessarily ungrammatical? Give the requested form and explain the distinction.
% mdparagraph:P268:end

\end{minipage}\par\addvspace{4pt}

\reserve{7\baselineskip}
% mdparagraph:P269:begin
\subsection*{Dependent questions}
% mdparagraph:P269:end

\reserve{5\baselineskip}
\par\begin{minipage}{\linewidth}\raggedright
% mdparagraph:P270:begin
\textbf{Q01.}
% mdparagraph:P270:end

% item: NEW-U09-Q01
% mdparagraph:P271:begin
Translate both clauses and explain the case of the question word. (a) Ic wāt, hwā ðone fisc hæbbe. (b) Ic wāt, hwǣm ðū ðone fisc sęlle.
% mdparagraph:P271:end

\end{minipage}\par\addvspace{4pt}

\par\begin{minipage}{\linewidth}\raggedright
% mdparagraph:P272:begin
\textbf{Q02.}
% mdparagraph:P272:end

% item: NEW-U09-Q02
% mdparagraph:P273:begin
Turn Hwæt hæfst ðū? into a dependent question after Saga mē. Keep the same addressee and use the present subjunctive.
% mdparagraph:P273:end

\end{minipage}\par\addvspace{4pt}

\par\begin{minipage}{\linewidth}\raggedright
% mdparagraph:P274:begin
\textbf{Q03.}
% mdparagraph:P274:end

% item: NEW-U09-Q03
% mdparagraph:P275:begin
Turn Hæfð se mann twēgen fiscas? into a dependent yes/\allowbreak{}no question after Saga mē. Use hwæðer and the present subjunctive.
% mdparagraph:P275:end

\end{minipage}\par\addvspace{4pt}

\par\begin{minipage}{\linewidth}\raggedright
% mdparagraph:P276:begin
\textbf{Q04.}
% mdparagraph:P276:end

% item: NEW-U09-Q04
% mdparagraph:P277:begin
Translate and distinguish the two uses of hwæðer. In (a), two men are under discussion. (a) Hwæðer mann is ðīn gefēra? (b) Saga mē, hwæðer se mann ðīn gefēra sīe.
% mdparagraph:P277:end

\end{minipage}\par\addvspace{4pt}

\par\begin{minipage}{\linewidth}\raggedright
% mdparagraph:P278:begin
\textbf{Q05.}
% mdparagraph:P278:end

% item: NEW-U09-Q05
% mdparagraph:P279:begin
Translate: I know whom you (sing.) help. Use a dependent question with hwǣm and the present subjunctive.
% mdparagraph:P279:end

\end{minipage}\par\addvspace{4pt}

\reserve{7\baselineskip}
% mdparagraph:P280:begin
\subsection*{Infinitives and object predicates}
% mdparagraph:P280:end

\reserve{5\baselineskip}
\par\begin{minipage}{\linewidth}\raggedright
% mdparagraph:P281:begin
\textbf{I01.}
% mdparagraph:P281:end

% item: NEW-U09-I01
% mdparagraph:P282:begin
Translate both sentences and identify who carries the fish. (a) Ic wille ðone fisc beran. (b) Ic hāte hine ðone fisc beran. Explain why beran has no person ending in either sentence.
% mdparagraph:P282:end

\end{minipage}\par\addvspace{4pt}

\par\begin{minipage}{\linewidth}\raggedright
% mdparagraph:P283:begin
\textbf{I02.}
% mdparagraph:P283:end

% item: NEW-U09-I02
% mdparagraph:P284:begin
The fisherman told two companions, “Berað ðone fisc.” Report this as a past command using hātan, hīe for the companions, and a plain infinitive.
% mdparagraph:P284:end

\end{minipage}\par\addvspace{4pt}

\par\begin{minipage}{\linewidth}\raggedright
% mdparagraph:P285:begin
\textbf{I03.}
% mdparagraph:P285:end

% item: NEW-U09-I03
% mdparagraph:P286:begin
Translate: I see the men carrying the fish (sing.). Use gesēon with an accusative actor and a plain infinitive; label both accusative noun phrases.
% mdparagraph:P286:end

\end{minipage}\par\addvspace{4pt}

\par\begin{minipage}{\linewidth}\raggedright
% mdparagraph:P287:begin
\textbf{I04.}
% mdparagraph:P287:end

% item: NEW-U09-I04
% mdparagraph:P288:begin
Translate both sentences using gewyrcean followed by an object and an agreeing predicate adjective: God made the man good. God made the men good. Explain the adjective endings.
% mdparagraph:P288:end

\end{minipage}\par\addvspace{4pt}

\par\begin{minipage}{\linewidth}\raggedright
% mdparagraph:P289:begin
\textbf{I05.}
% mdparagraph:P289:end

% item: NEW-U09-I05
% mdparagraph:P290:begin
Translate and distinguish the two complements: Hē hēt ðone mann cyning. Hē hēt ðone mann gān. State the case/\allowbreak{}function of cyning and the form/\allowbreak{}function of gān.
% mdparagraph:P290:end

\end{minipage}\par\addvspace{4pt}

\reserve{7\baselineskip}
% mdparagraph:P291:begin
\subsection*{Unreal wishes}
% mdparagraph:P291:end

\reserve{5\baselineskip}
\par\begin{minipage}{\linewidth}\raggedright
% mdparagraph:P292:begin
\textbf{W01.}
% mdparagraph:P292:end

% item: NEW-U09-W01
% mdparagraph:P293:begin
You have no ship. Express the present unreal wish “Would that I had one ship!” using the preterite subjunctive of habban.
% mdparagraph:P293:end

\end{minipage}\par\addvspace{4pt}

\par\begin{minipage}{\linewidth}\raggedright
% mdparagraph:P294:begin
\textbf{W02.}
% mdparagraph:P294:end

% item: NEW-U09-W02
% mdparagraph:P295:begin
Translate these in their stated settings: (a) a report about last year, Ic hæfde ān scip.; (b) a wish by someone who has no ship now, Hæfde ic ān scip! Explain what the shared spelling hæfde can and cannot show.
% mdparagraph:P295:end

\end{minipage}\par\addvspace{4pt}

\reserve{7\baselineskip}
% mdparagraph:P296:begin
\section*{10. The fisherman and his companion}
% mdparagraph:P296:end

% mdparagraph:P297:begin
New teaching material. Read the connected exchange, then answer the questions below.
% mdparagraph:P297:end

% composition: NEW-U09-READING-01
% mdparagraph:P298:begin
\par\begin{minipage}{\linewidth}\raggedright
\readingtext{Se fiscere and his gefēra standað on ðǣm strande. Se gefēra cwið: “Saga mē, hwæt ðū hæbbe.” Se fiscere sægð, ðæt hē męte hæbbe. Hē hǣtt his gefēran ðone fisc beran. Se gefēra birð ðone fisc. Se fiscere cwið: “Ic gesēo ðē ðone fisc beran. Nū wāt ic, ðæt ðū mē hilpst.”}
\end{minipage}\par
% mdparagraph:P298:end

\reserve{7\baselineskip}
% mdparagraph:P299:begin
\subsection*{Reading practice}
% mdparagraph:P299:end

\reserve{5\baselineskip}
\par\begin{minipage}{\linewidth}\raggedright
% mdparagraph:P300:begin
\textbf{R01.}
% mdparagraph:P300:end

% item: NEW-U09-R01
% mdparagraph:P301:begin
In the controlled reading, who requests information, and who is reported to have food? Answer in English and cite the two relevant clauses.
% mdparagraph:P301:end

\end{minipage}\par\addvspace{4pt}

\par\begin{minipage}{\linewidth}\raggedright
% mdparagraph:P302:begin
\textbf{R02.}
% mdparagraph:P302:end

% item: NEW-U09-R02
% mdparagraph:P303:begin
In the controlled reading, compare hwæt ðū hæbbe, ðæt hē męte hæbbe, and ðæt ðū mē hilpst. Identify each clause type and explain the selected moods.
% mdparagraph:P303:end

\end{minipage}\par\addvspace{4pt}

\par\begin{minipage}{\linewidth}\raggedright
% mdparagraph:P304:begin
\textbf{R03.}
% mdparagraph:P304:end

% item: NEW-U09-R03
% mdparagraph:P305:begin
Identify the commanded actor in Hē hǣtt his gefēran ðone fisc beran and the observed actor in Ic gesēo ðē ðone fisc beran. For each sentence, distinguish that actor from the fish and identify the finite verb.
% mdparagraph:P305:end

\end{minipage}\par\addvspace{4pt}

\par\begin{minipage}{\linewidth}\raggedright
% mdparagraph:P306:begin
\textbf{R04.}
% mdparagraph:P306:end

% item: NEW-U09-R04
% mdparagraph:P307:begin
Translate the complete controlled reading, preserving the speakers, the command, the carrying action and the final distinction between report and asserted knowledge.
% mdparagraph:P307:end

\end{minipage}\par\addvspace{4pt}

\reserve{7\baselineskip}
% mdparagraph:P308:begin
\section*{11. The Choice of Occupations}
% mdparagraph:P308:end

% mdparagraph:P309:begin
Exact source text from Cook’s \emph{A First Book in Old English}. The Teacher asks the Counsellor to rank occupations. Several expressions are supplied in the local helps for reading; their unfamiliar forms are not required for independent composition.
% mdparagraph:P309:end

\par\begin{minipage}{\linewidth}\raggedright
% source: FB-II-CHOICE-OCCUPATIONS-01
% mdparagraph:P310:begin
\readingtext{\textbf{Teacher.} Hwæt sægst ðū, wīsa? Hwilc cræft ðē is geðūht betweox ðās furðra wesan?}
% mdparagraph:P310:end

% mdparagraph:P311:begin
\readingtext{\textbf{Counsellor.} Ic sęcge ðē, mē is geðūht Godes ðēowdōm betweoh ðās cræftas ealdorscipe healdan, swā-swā hit is gerǣd on godspelle, “Fyrmest sēceað rīce Godes, and rihtwīsnesse his, and ðās ðing eall bēoð tōgēiecte ēow.”}
% mdparagraph:P311:end

% mdparagraph:P312:begin
\readingtext{\textbf{Teacher.} And hwilc ðē is geðūht betweox woruldcræftas healdan ealdordōm?}
% mdparagraph:P312:end

% mdparagraph:P313:begin
\readingtext{\textbf{Counsellor.} Eorðtilð, for-ðām se ierðling ūs ealle fēt.}
% mdparagraph:P313:end

\end{minipage}\par
\reserve{7\baselineskip}
% mdparagraph:P314:begin
\subsection*{Local help}
% mdparagraph:P314:end

% teaching: NEW-U09-SOURCE-HELPS
% teaching: NEW-U09-GRAM-u09-source-seceath
% teaching: NEW-U09-VOC-thyncan
% teaching: NEW-U09-VOC-betweox
% teaching: NEW-U09-VOC-furthra
% teaching: NEW-U09-VOC-swa-swa
% teaching: NEW-U09-VOC-fyrmest
% teaching: NEW-U09-VOC-togeiecan
% teaching: NEW-U09-VOC-woruldcraeft
% teaching: NEW-U09-VOC-ealdordom
% teaching: NEW-U09-VOC-fedan
% support: U09-SRC-HELP01
% mdparagraph:P315:begin
\textbf{1. wīsa.} “Wise man; wise one”, in direct address. It is the weak masculine singular of wīs used as a noun. This personal use is given here for reading.
% mdparagraph:P315:end

% support: U09-SRC-HELP02
% mdparagraph:P316:begin
\textbf{2. ðē is geðūht … wesan /\allowbreak{} mē is geðūht … healdan.} “To you it seems … to be /\allowbreak{} to me it seems … to hold.” Read these as complete expressions. The person to whom it seems is dative. Cook identifies furðra as nominative singular; these source expressions do not follow one uniform accusative-actor pattern like commanding hātan.
% mdparagraph:P316:end

% support: U09-SRC-HELP03
% mdparagraph:P317:begin
\textbf{3. betweox /\allowbreak{} betweoh ðās (cræftas).} “Among these (occupations).” The two spellings mean the same thing here. The following phrase is accusative plural. In the first question, cræftas is understood after ðās.
% mdparagraph:P317:end

% support: U09-SRC-HELP04
% mdparagraph:P318:begin
\textbf{4. furðra.} “Foremost; prior.” This is a weak comparative, masculine nominative singular, as Cook’s note 14 states. This one form and sense are given for the passage.
% mdparagraph:P318:end

% support: U09-SRC-HELP05
% mdparagraph:P319:begin
\textbf{5. swā-swā hit is gerǣd on godspelle.} “Just as it is read in the Gospel.” Swā-swā means “just as”; gerǣd is the past participle in this passive expression. Read the whole phrase with this help.
% mdparagraph:P319:end

% support: U09-SRC-HELP06
% mdparagraph:P320:begin
\textbf{6. Fyrmest sēceað.} “First seek”, a plural command. Sēceað is the form printed in the source; recognize it beside the selected plural command sēcað.
% mdparagraph:P320:end

% support: U09-SRC-HELP07
% mdparagraph:P321:begin
\textbf{7. ðās ðing eall bēoð tōgēiecte ēow.} “All these things will be added to you.” Eall agrees with neuter plural ðing. Bēoð is third-person plural here; ēow is the dative recipient. Tōgēiecte is a predicate participle. Keep its printed -e: Cook’s §62 describes this predicate use, so it should not be silently changed to an expected neuter ending.
% mdparagraph:P321:end

% support: U09-SRC-HELP08
% mdparagraph:P322:begin
\textbf{8. woruldcræftas /\allowbreak{} ealdordōm.} “Worldly occupations” (accusative plural) /\allowbreak{} “chief place, primacy” (accusative singular). Cook prints woruldcræftas but records the manuscript words cræftas woruld in note 20.
% mdparagraph:P322:end

% support: U09-SRC-HELP09
% mdparagraph:P323:begin
\textbf{9. ūs ealle fēt.} “Feeds us all.” Ūs is accusative plural; ealle refers to the people fed and is masculine plural. Fēt is present indicative third singular of fēdan. This exact form is given for reading; the prefixed āfēdan table does not supply a complete system for fēdan.
% mdparagraph:P323:end

% support: U09-SRC-HELP10
% mdparagraph:P324:begin
\textbf{10. the three conjectural is insertions.} Cook calls the is in each Teacher question and in swā-swā hit is gerǣd conjectural: they are absent from his manuscripts. The is in mē is geðūht instead belongs to note 15. The printed text here retains Cook’s three insertions.
% mdparagraph:P324:end

% support: U09-SRC-HELP11
% mdparagraph:P325:begin
\textbf{11. Godes ðēowdōm; rihtwīsnesse his.} “God’s service, service of God”; “his righteousness”. The genitives identify the divine possessor or the one served. His follows its noun in the second phrase.
% mdparagraph:P325:end

\reserve{7\baselineskip}
% mdparagraph:P326:begin
\subsection*{Cook’s original notes}
% mdparagraph:P326:end

% mdparagraph:P327:begin
The numbers are Cook’s original note numbers. “Lat.” identifies his Latin comparison; “MSS.” means manuscripts. His section references point to the grammar in his \emph{First Book}.
% mdparagraph:P327:end

% mdtable:T73:begin
\par\begingroup
\fontsize{10.5}{13.2}\selectfont
\setlength{\LTleft}{0pt}\setlength{\LTright}{\fill}
\noindent
\begin{tabular}{@{}>{\raggedright\arraybackslash}p{22.72mm}>{\raggedright\arraybackslash}p{38.74mm}>{\raggedright\arraybackslash}p{56.92mm}@{}}
\toprule
\textbf{Note} & \textbf{Anchor} & \textbf{Original note} \\
\midrule
11 & cræft & Lat. ars. \\
12 & is (three occurrences) & Conjectural; not in the MSS. \\
13 & geðūht & See ðyncan. \\
14 & furðra & Lat. prior. Nom. sing. \\
15 & mē is geðūht & Lat. videtur. \\
16 & ðēowdōm & See 143 and 149. \\
17 & ealdorscipe; ealdordōm & Lat. primatum. \\
18 & rihtwīsnesse & See 144. \\
19 & tōgēiecte & Lat. adjicientur. See tōgēiecan, and 62. \\
20 & woruldcræftas & Lat. artes seculares. MS. cræftas woruld. \\
21 & Eorðtilð & Lat. agricultura. See 147. \\
22 & ierðling & Lat. arator. \\
23 & fēt & See fēdan. \\
\bottomrule
\end{tabular}
\par\endgroup

% mdtable:T73:end
\reserve{7\baselineskip}
% mdparagraph:P328:begin
\subsection*{Supported source practice}
% mdparagraph:P328:end

% mdparagraph:P329:begin
Use the text, vocabulary and local helps. These tasks ask you to understand the supported source; they do not require unaided recall of the unfamiliar source forms.
% mdparagraph:P329:end

\reserve{5\baselineskip}
\par\begin{minipage}{\linewidth}\raggedright
% mdparagraph:P330:begin
\textbf{S01.}
% mdparagraph:P330:end

% item: NEW-U09-S01
% mdparagraph:P331:begin
In the source dialogue, which occupation receives first place overall, and which receives first place among worldly occupations? Explain why the two answers do not contradict one another.
% mdparagraph:P331:end

\end{minipage}\par\addvspace{4pt}

\par\begin{minipage}{\linewidth}\raggedright
% mdparagraph:P332:begin
\textbf{S02.}
% mdparagraph:P332:end

% item: NEW-U09-S02
% mdparagraph:P333:begin
Identify the beginning and end of the Gospel quotation in the source dialogue. What does it command, and what does it promise? Use the printed helps for unfamiliar forms.
% mdparagraph:P333:end

\end{minipage}\par\addvspace{4pt}

\par\begin{minipage}{\linewidth}\raggedright
% mdparagraph:P334:begin
\textbf{S03.}
% mdparagraph:P334:end

% item: NEW-U09-S03
% mdparagraph:P335:begin
In Ic sęcge ðē, mē is geðūht, explain the different roles of ðē and mē using the local help. Who is addressed, and whose judgment follows?
% mdparagraph:P335:end

\end{minipage}\par\addvspace{4pt}

\par\begin{minipage}{\linewidth}\raggedright
% mdparagraph:P336:begin
\textbf{S04.}
% mdparagraph:P336:end

% item: NEW-U09-S04
% mdparagraph:P337:begin
What reason does the Counsellor give for agriculture's place? Explain the reference of ūs ealle in the final turn and identify the subject of fēt.
% mdparagraph:P337:end

\end{minipage}\par\addvspace{4pt}

\par\begin{minipage}{\linewidth}\raggedright
% mdparagraph:P338:begin
\textbf{S05.}
% mdparagraph:P338:end

% item: NEW-U09-S05
% mdparagraph:P339:begin
Translate all four source turns, including the complete Gospel quotation. Preserve the distinction between the overall ranking and the later question about worldly occupations.
% mdparagraph:P339:end

\end{minipage}\par\addvspace{4pt}

\reserve{7\baselineskip}
% mdparagraph:P340:begin
\section*{12. Cook II.9–12}
% mdparagraph:P340:end

% mdparagraph:P341:begin
Translate the complete English prompts into Old English. Cook’s numbering, parentheses, brackets and grammar references are retained. “Sing.” and “plur.” specify number; “om.” asks you to omit the English word indicated. Section numbers refer to Cook’s grammar. Use the local help printed with the three supported expressions.
% mdparagraph:P341:end

% mdparagraph:P342:begin
In a “let us” or “let …” exhortation, use the appropriate present subjunctive unless another construction is expressly requested. Direct address uses the nominative-shaped noun and a weak accompanying adjective. In the reported-speech item, follow the prompt as printed; a content conjunction has not been silently added.
% mdparagraph:P342:end

\reserve{7\baselineskip}
% mdparagraph:P343:begin
\subsection*{Exercise 9}
% mdparagraph:P343:end

\reserve{5\baselineskip}
\par\begin{minipage}{\linewidth}\raggedright
% mdparagraph:P344:begin
\textbf{Cook II.9.1.}
% mdparagraph:P344:end

% item: EX-II-09-01
% mdparagraph:P345:begin
Sparks fly.
% mdparagraph:P345:end

\end{minipage}\par\addvspace{4pt}

\par\begin{minipage}{\linewidth}\raggedright
% mdparagraph:P346:begin
\textbf{Cook II.9.2.}
% mdparagraph:P346:end

% item: EX-II-09-02
% mdparagraph:P347:begin
The house is empty.
% mdparagraph:P347:end

\end{minipage}\par\addvspace{4pt}

\par\begin{minipage}{\linewidth}\raggedright
% mdparagraph:P348:begin
\textbf{Cook II.9.3.}
% mdparagraph:P348:end

% item: EX-II-09-03
% mdparagraph:P349:begin
The ships are empty.
% mdparagraph:P349:end

\end{minipage}\par\addvspace{4pt}

\par\begin{minipage}{\linewidth}\raggedright
% mdparagraph:P350:begin
\textbf{Cook II.9.4.}
% mdparagraph:P350:end

% item: EX-II-09-04
% mdparagraph:P351:begin
The merchant sells needles, awls, fish-hooks, and plowshares.
% mdparagraph:P351:end

\end{minipage}\par\addvspace{4pt}

\par\begin{minipage}{\linewidth}\raggedright
% mdparagraph:P352:begin
\textbf{Cook II.9.5.}
% mdparagraph:P352:end

% item: EX-II-09-05
% mdparagraph:P353:begin
The blacksmith makes plowshares in his smithy.
% mdparagraph:P353:end

\end{minipage}\par\addvspace{4pt}

\par\begin{minipage}{\linewidth}\raggedright
% mdparagraph:P354:begin
\textbf{Cook II.9.6.}
% mdparagraph:P354:end

% item: EX-II-09-06
% mdparagraph:P355:begin
The oil and the wine are good.
% mdparagraph:P355:end

\end{minipage}\par\addvspace{4pt}

\par\begin{minipage}{\linewidth}\raggedright
% mdparagraph:P356:begin
\textbf{Cook II.9.7.}
% mdparagraph:P356:end

% item: EX-II-09-07
% mdparagraph:P357:begin
The fisherman rows over the sea.
% mdparagraph:P357:end

\end{minipage}\par\addvspace{4pt}

\par\begin{minipage}{\linewidth}\raggedright
% mdparagraph:P358:begin
\textbf{Cook II.9.8.}
% mdparagraph:P358:end

% item: EX-II-09-08
% mdparagraph:P359:begin
The carpenter supports his son.
% mdparagraph:P359:end

\end{minipage}\par\addvspace{4pt}

\par\begin{minipage}{\linewidth}\raggedright
% mdparagraph:P360:begin
\textbf{Cook II.9.9.}
% mdparagraph:P360:end

% item: EX-II-09-09
% mdparagraph:P361:begin
The tailor bought the needle.
% mdparagraph:P361:end

\end{minipage}\par\addvspace{4pt}

\par\begin{minipage}{\linewidth}\raggedright
% mdparagraph:P362:begin
\textbf{Cook II.9.10.}
% mdparagraph:P362:end

% item: EX-II-09-10
% mdparagraph:P363:begin
The bellows blow.
% mdparagraph:P363:end

\end{minipage}\par\addvspace{4pt}

\par\begin{minipage}{\linewidth}\raggedright
% mdparagraph:P364:begin
\textbf{Cook II.9.11.}
% mdparagraph:P364:end

% item: EX-II-09-11
% mdparagraph:P365:begin
The hammer smites,
% mdparagraph:P365:end

\end{minipage}\par\addvspace{4pt}

\par\begin{minipage}{\linewidth}\raggedright
% mdparagraph:P366:begin
\textbf{Cook II.9.12.}
% mdparagraph:P366:end

% item: EX-II-09-12
% mdparagraph:P367:begin
Let us seek the kingdom of God.
% mdparagraph:P367:end

\end{minipage}\par\addvspace{4pt}

\par\begin{minipage}{\linewidth}\raggedright
% mdparagraph:P368:begin
\textbf{Cook II.9.13.}
% mdparagraph:P368:end

% item: EX-II-09-13
% mdparagraph:P369:begin
He will not live with the blacksmith.
% mdparagraph:P369:end

\end{minipage}\par\addvspace{4pt}

\par\begin{minipage}{\linewidth}\raggedright
% mdparagraph:P370:begin
\textbf{Cook II.9.14.}
% mdparagraph:P370:end

% item: EX-II-09-14
% mdparagraph:P371:begin
We had rather live [To us it is preferable to sojourn] with the farmer than with the tailor.
% mdparagraph:P371:end

\end{minipage}\par\addvspace{4pt}

\par\begin{minipage}{\linewidth}\raggedright
% mdparagraph:P372:begin
\textbf{Cook II.9.15.}
% mdparagraph:P372:end

% item: EX-II-09-15
% mdparagraph:P373:begin
Gold is better (66) than tin.
% mdparagraph:P373:end

\end{minipage}\par\addvspace{4pt}

\par\begin{minipage}{\linewidth}\raggedright
% mdparagraph:P374:begin
\textbf{Cook II.9.16.}
% mdparagraph:P374:end

% item: EX-II-09-16
% mdparagraph:P375:begin
The carpenter makes a hole.
% mdparagraph:P375:end

\end{minipage}\par\addvspace{4pt}

\reserve{7\baselineskip}
% mdparagraph:P376:begin
\subsection*{Exercise 10}
% mdparagraph:P376:end

\reserve{5\baselineskip}
\par\begin{minipage}{\linewidth}\raggedright
% mdparagraph:P377:begin
\textbf{Cook II.10.1.}
% mdparagraph:P377:end

% item: EX-II-10-01
% mdparagraph:P378:begin
The tree is useful to the carpenter.
% mdparagraph:P378:end

\end{minipage}\par\addvspace{4pt}

\par\begin{minipage}{\linewidth}\raggedright
% mdparagraph:P379:begin
\textbf{Cook II.10.2.}
% mdparagraph:P379:end

% item: EX-II-10-02
% mdparagraph:P380:begin
Our horses have fodder.
% mdparagraph:P380:end

\end{minipage}\par\addvspace{4pt}

\par\begin{minipage}{\linewidth}\raggedright
% mdparagraph:P381:begin
\textbf{Cook II.10.3.}
% mdparagraph:P381:end

% item: EX-II-10-03
% mdparagraph:P382:begin
O priest, read the Gospel.
% mdparagraph:P382:end

\end{minipage}\par\addvspace{4pt}

\par\begin{minipage}{\linewidth}\raggedright
% mdparagraph:P383:begin
\textbf{Cook II.10.4.}
% mdparagraph:P383:end

% item: EX-II-10-04
% mdparagraph:P384:begin
Buy (sing.) wine and oil.
% mdparagraph:P384:end

\end{minipage}\par\addvspace{4pt}

\par\begin{minipage}{\linewidth}\raggedright
% mdparagraph:P385:begin
\textbf{Cook II.10.5.}
% mdparagraph:P385:end

% item: EX-II-10-05
% mdparagraph:P386:begin
O good fellows, row over the waters.
% mdparagraph:P386:end

\end{minipage}\par\addvspace{4pt}

\par\begin{minipage}{\linewidth}\raggedright
% mdparagraph:P387:begin
\textbf{Cook II.10.6.}
% mdparagraph:P387:end

% item: EX-II-10-06
% mdparagraph:P388:begin
Her son is rich.
% mdparagraph:P388:end

\end{minipage}\par\addvspace{4pt}

\par\begin{minipage}{\linewidth}\raggedright
% mdparagraph:P389:begin
\textbf{Cook II.10.7.}
% mdparagraph:P389:end

% item: EX-II-10-07
% mdparagraph:P390:begin
Ply (sing.) your trade.
% mdparagraph:P390:end

\end{minipage}\par\addvspace{4pt}

\par\begin{minipage}{\linewidth}\raggedright
% mdparagraph:P391:begin
\textbf{Cook II.10.8.}
% mdparagraph:P391:end

% item: EX-II-10-08
% mdparagraph:P392:begin
God gave us his Gospel.
% mdparagraph:P392:end

\end{minipage}\par\addvspace{4pt}

\par\begin{minipage}{\linewidth}\raggedright
% mdparagraph:P393:begin
\textbf{Cook II.10.9.}
% mdparagraph:P393:end

% item: EX-II-10-09
% mdparagraph:P394:begin
The farmer practises agriculture.
% mdparagraph:P394:end

\end{minipage}\par\addvspace{4pt}

\par\begin{minipage}{\linewidth}\raggedright
% mdparagraph:P395:begin
\textbf{Cook II.10.10.}
% mdparagraph:P395:end

% item: EX-II-10-10
% mdparagraph:P396:begin
The service of God [God’s service] is better (66) than other occupations.
% mdparagraph:P396:end

\end{minipage}\par\addvspace{4pt}

\par\begin{minipage}{\linewidth}\raggedright
% mdparagraph:P397:begin
\textbf{Cook II.10.11.}
% mdparagraph:P397:end

% item: EX-II-10-11
% mdparagraph:P398:begin
Which of you holds (om. the) chief place?
% mdparagraph:P398:end

\end{minipage}\par\addvspace{4pt}

\par\begin{minipage}{\linewidth}\raggedright
% mdparagraph:P399:begin
\textbf{Cook II.10.12.}
% mdparagraph:P399:end

% item: EX-II-10-12
% mdparagraph:P400:begin
My gain is little.
% mdparagraph:P400:end

\end{minipage}\par\addvspace{4pt}

\par\begin{minipage}{\linewidth}\raggedright
% mdparagraph:P401:begin
\textbf{Cook II.10.13.}
% mdparagraph:P401:end

% item: EX-II-10-13
% mdparagraph:P402:begin
My labor is great.
% mdparagraph:P402:end

\end{minipage}\par\addvspace{4pt}

% teaching: NEW-U09-COOK-HELP-IVORY
% mdparagraph:P403:begin
\textbf{Local help for II.10.14:} \textbf{elpes bān} — “ivory”, literally “elephant’s bone”. Elpes is genitive singular; bān is neuter accusative singular as the thing sold. Use this whole expression as given.
% mdparagraph:P403:end

\reserve{5\baselineskip}
\par\begin{minipage}{\linewidth}\raggedright
% mdparagraph:P404:begin
\textbf{Cook II.10.14.}
% mdparagraph:P404:end

% item: EX-II-10-14
% mdparagraph:P405:begin
I sold ivory and silk.
% mdparagraph:P405:end

\end{minipage}\par\addvspace{4pt}

\par\begin{minipage}{\linewidth}\raggedright
% mdparagraph:P406:begin
\textbf{Cook II.10.15.}
% mdparagraph:P406:end

% item: EX-II-10-15
% mdparagraph:P407:begin
O, were I king of this people!
% mdparagraph:P407:end

\end{minipage}\par\addvspace{4pt}

\reserve{7\baselineskip}
% mdparagraph:P408:begin
\subsection*{Exercise 11}
% mdparagraph:P408:end

\reserve{5\baselineskip}
\par\begin{minipage}{\linewidth}\raggedright
% mdparagraph:P409:begin
\textbf{Cook II.11.1.}
% mdparagraph:P409:end

% item: EX-II-11-01
% mdparagraph:P410:begin
The monks have food and drink.
% mdparagraph:P410:end

\end{minipage}\par\addvspace{4pt}

\par\begin{minipage}{\linewidth}\raggedright
% mdparagraph:P411:begin
\textbf{Cook II.11.2.}
% mdparagraph:P411:end

% item: EX-II-11-02
% mdparagraph:P412:begin
Settle [quell] your disputes.
% mdparagraph:P412:end

\end{minipage}\par\addvspace{4pt}

\par\begin{minipage}{\linewidth}\raggedright
% mdparagraph:P413:begin
\textbf{Cook II.11.3.}
% mdparagraph:P413:end

% item: EX-II-11-03
% mdparagraph:P414:begin
He suffered shipwreck.
% mdparagraph:P414:end

\end{minipage}\par\addvspace{4pt}

\par\begin{minipage}{\linewidth}\raggedright
% mdparagraph:P415:begin
\textbf{Cook II.11.4.}
% mdparagraph:P415:end

% item: EX-II-11-04
% mdparagraph:P416:begin
He abandoned his wife.
% mdparagraph:P416:end

\end{minipage}\par\addvspace{4pt}

\par\begin{minipage}{\linewidth}\raggedright
% mdparagraph:P417:begin
\textbf{Cook II.11.5.}
% mdparagraph:P417:end

% item: EX-II-11-05
% mdparagraph:P418:begin
Be what [that] you (sing.) ought to be.
% mdparagraph:P418:end

\end{minipage}\par\addvspace{4pt}

\par\begin{minipage}{\linewidth}\raggedright
% mdparagraph:P419:begin
\textbf{Cook II.11.6.}
% mdparagraph:P419:end

% item: EX-II-11-06
% mdparagraph:P420:begin
They do not wish to be what they ought to be.
% mdparagraph:P420:end

\end{minipage}\par\addvspace{4pt}

\par\begin{minipage}{\linewidth}\raggedright
% mdparagraph:P421:begin
\textbf{Cook II.11.7.}
% mdparagraph:P421:end

% item: EX-II-11-07
% mdparagraph:P422:begin
He went aboard the ship with his wares.
% mdparagraph:P422:end

\end{minipage}\par\addvspace{4pt}

\par\begin{minipage}{\linewidth}\raggedright
% mdparagraph:P423:begin
\textbf{Cook II.11.8.}
% mdparagraph:P423:end

% item: EX-II-11-08
% mdparagraph:P424:begin
O good workmen, ply your trades.
% mdparagraph:P424:end

\end{minipage}\par\addvspace{4pt}

\par\begin{minipage}{\linewidth}\raggedright
% mdparagraph:P425:begin
\textbf{Cook II.11.9.}
% mdparagraph:P425:end

% item: EX-II-11-09
% mdparagraph:P426:begin
I tell you (sing.), the shoemaker buys awls, and the fisherman buys (fish-) hooks.
% mdparagraph:P426:end

\end{minipage}\par\addvspace{4pt}

\par\begin{minipage}{\linewidth}\raggedright
% mdparagraph:P427:begin
\textbf{Cook II.11.10.}
% mdparagraph:P427:end

% item: EX-II-11-10
% mdparagraph:P428:begin
Sometimes he reads, because he is wise.
% mdparagraph:P428:end

\end{minipage}\par\addvspace{4pt}

\par\begin{minipage}{\linewidth}\raggedright
% mdparagraph:P429:begin
\textbf{Cook II.11.11.}
% mdparagraph:P429:end

% item: EX-II-11-11
% mdparagraph:P430:begin
The king buys horses.
% mdparagraph:P430:end

\end{minipage}\par\addvspace{4pt}

\par\begin{minipage}{\linewidth}\raggedright
% mdparagraph:P431:begin
\textbf{Cook II.11.12.}
% mdparagraph:P431:end

% item: EX-II-11-12
% mdparagraph:P432:begin
Righteousness is better than gold.
% mdparagraph:P432:end

\end{minipage}\par\addvspace{4pt}

\par\begin{minipage}{\linewidth}\raggedright
% mdparagraph:P433:begin
\textbf{Cook II.11.13.}
% mdparagraph:P433:end

% item: EX-II-11-13
% mdparagraph:P434:begin
They struck me.
% mdparagraph:P434:end

\end{minipage}\par\addvspace{4pt}

\par\begin{minipage}{\linewidth}\raggedright
% mdparagraph:P435:begin
\textbf{Cook II.11.14.}
% mdparagraph:P435:end

% item: EX-II-11-14
% mdparagraph:P436:begin
The farmer beats him.
% mdparagraph:P436:end

\end{minipage}\par\addvspace{4pt}

\par\begin{minipage}{\linewidth}\raggedright
% mdparagraph:P437:begin
\textbf{Cook II.11.15.}
% mdparagraph:P437:end

% item: EX-II-11-15
% mdparagraph:P438:begin
We suffer the loss of copper and tin.
% mdparagraph:P438:end

\end{minipage}\par\addvspace{4pt}

\par\begin{minipage}{\linewidth}\raggedright
% mdparagraph:P439:begin
\textbf{Cook II.11.16.}
% mdparagraph:P439:end

% item: EX-II-11-16
% mdparagraph:P440:begin
The princes of the earth go aboard their ships.
% mdparagraph:P440:end

\end{minipage}\par\addvspace{4pt}

\par\begin{minipage}{\linewidth}\raggedright
% mdparagraph:P441:begin
\textbf{Cook II.11.17.}
% mdparagraph:P441:end

% item: EX-II-11-17
% mdparagraph:P442:begin
The priest benefits (164. e) the laymen.
% mdparagraph:P442:end

\end{minipage}\par\addvspace{4pt}

\reserve{7\baselineskip}
% mdparagraph:P443:begin
\subsection*{Exercise 12}
% mdparagraph:P443:end

\reserve{5\baselineskip}
\par\begin{minipage}{\linewidth}\raggedright
% mdparagraph:P444:begin
\textbf{Cook II.12.1.}
% mdparagraph:P444:end

% item: EX-II-12-01
% mdparagraph:P445:begin
The counsellor is wiser (64) than the warrior.
% mdparagraph:P445:end

\end{minipage}\par\addvspace{4pt}

\par\begin{minipage}{\linewidth}\raggedright
% mdparagraph:P446:begin
\textbf{Cook II.12.2.}
% mdparagraph:P446:end

% item: EX-II-12-02
% mdparagraph:P447:begin
The elephant is larger than the horse.
% mdparagraph:P447:end

\end{minipage}\par\addvspace{4pt}

\par\begin{minipage}{\linewidth}\raggedright
% mdparagraph:P448:begin
\textbf{Cook II.12.3.}
% mdparagraph:P448:end

% item: EX-II-12-03
% mdparagraph:P449:begin
God created all (om. the) fishes in (om. the) sea.
% mdparagraph:P449:end

\end{minipage}\par\addvspace{4pt}

\par\begin{minipage}{\linewidth}\raggedright
% mdparagraph:P450:begin
\textbf{Cook II.12.4.}
% mdparagraph:P450:end

% item: EX-II-12-04
% mdparagraph:P451:begin
The cattle seek fodder.
% mdparagraph:P451:end

\end{minipage}\par\addvspace{4pt}

\par\begin{minipage}{\linewidth}\raggedright
% mdparagraph:P452:begin
\textbf{Cook II.12.5.}
% mdparagraph:P452:end

% item: EX-II-12-05
% mdparagraph:P453:begin
God commanded the earth to bring forth grass.
% mdparagraph:P453:end

\end{minipage}\par\addvspace{4pt}

\par\begin{minipage}{\linewidth}\raggedright
% mdparagraph:P454:begin
\textbf{Cook II.12.6.}
% mdparagraph:P454:end

% item: EX-II-12-06
% mdparagraph:P455:begin
On the third day God gathered the waters that were under the heavens.
% mdparagraph:P455:end

\end{minipage}\par\addvspace{4pt}

% teaching: NEW-U09-COOK-HELP-LIVING
% mdparagraph:P456:begin
\textbf{Local help for II.12.7:} \textbf{libbende gesceafta} — “living creatures”, feminine accusative plural. Use this whole phrase as given. Libbende uses the documented feminine -e participial ending; the earlier neuter libbendu would be a different agreement form.
% mdparagraph:P456:end

\reserve{5\baselineskip}
\par\begin{minipage}{\linewidth}\raggedright
% mdparagraph:P457:begin
\textbf{Cook II.12.7.}
% mdparagraph:P457:end

% item: EX-II-12-07
% mdparagraph:P458:begin
Let the earth bring forth living creatures (the verb precedes).
% mdparagraph:P458:end

\end{minipage}\par\addvspace{4pt}

\par\begin{minipage}{\linewidth}\raggedright
% mdparagraph:P459:begin
\textbf{Cook II.12.8.}
% mdparagraph:P459:end

% item: EX-II-12-08
% mdparagraph:P460:begin
God made man good (173).
% mdparagraph:P460:end

\end{minipage}\par\addvspace{4pt}

\par\begin{minipage}{\linewidth}\raggedright
% mdparagraph:P461:begin
\textbf{Cook II.12.9.}
% mdparagraph:P461:end

% item: EX-II-12-09
% mdparagraph:P462:begin
He is (om. the) light that gives light to (ālīehtan) all things.
% mdparagraph:P462:end

\end{minipage}\par\addvspace{4pt}

\par\begin{minipage}{\linewidth}\raggedright
% mdparagraph:P463:begin
\textbf{Cook II.12.10.}
% mdparagraph:P463:end

% item: EX-II-12-10
% mdparagraph:P464:begin
Some of them (151) were rich, and some were wise.
% mdparagraph:P464:end

\end{minipage}\par\addvspace{4pt}

\par\begin{minipage}{\linewidth}\raggedright
% mdparagraph:P465:begin
\textbf{Cook II.12.11.}
% mdparagraph:P465:end

% item: EX-II-12-11
% mdparagraph:P466:begin
The night has many stars.
% mdparagraph:P466:end

\end{minipage}\par\addvspace{4pt}

\par\begin{minipage}{\linewidth}\raggedright
% mdparagraph:P467:begin
\textbf{Cook II.12.12.}
% mdparagraph:P467:end

% item: EX-II-12-12
% mdparagraph:P468:begin
The light shineth in (on) darkness (plur.).
% mdparagraph:P468:end

\end{minipage}\par\addvspace{4pt}

% teaching: NEW-U09-COOK-HELP-365
% teaching: NEW-U09-VOC-hund
% teaching: NEW-U09-VOC-siextig
% mdparagraph:P469:begin
\textbf{Local help for II.12.13:} \textbf{ðrēo hund and fīf and siextig daga} — “three hundred and sixty-five days”. Use this whole phrase as given. Hund is “hundred”, siextig “sixty”; ðrēo counts the neuter hundred-unit, and daga is genitive plural. The phrase is composed from Cook’s numeral rules, not quoted as a sentence from his reader. You are not being asked to recall a complete higher-numeral system.
% mdparagraph:P469:end

\reserve{5\baselineskip}
\par\begin{minipage}{\linewidth}\raggedright
% mdparagraph:P470:begin
\textbf{Cook II.12.13.}
% mdparagraph:P470:end

% item: EX-II-12-13
% mdparagraph:P471:begin
The year has three hundred and sixty-five (78) days.
% mdparagraph:P471:end

\end{minipage}\par\addvspace{4pt}

\par\begin{minipage}{\linewidth}\raggedright
% mdparagraph:P472:begin
\textbf{Cook II.12.14.}
% mdparagraph:P472:end

% item: EX-II-12-14
% mdparagraph:P473:begin
The carpenter has two sons.
% mdparagraph:P473:end

\end{minipage}\par\addvspace{4pt}

\par\begin{minipage}{\linewidth}\raggedright
% mdparagraph:P474:begin
\textbf{Cook II.12.15.}
% mdparagraph:P474:end

% item: EX-II-12-15
% mdparagraph:P475:begin
The soldier has eight houses.
% mdparagraph:P475:end

\end{minipage}\par\addvspace{4pt}

\par\begin{minipage}{\linewidth}\raggedright
% mdparagraph:P476:begin
\textbf{Cook II.12.16.}
% mdparagraph:P476:end

% item: EX-II-12-16
% mdparagraph:P477:begin
The peril was great.
% mdparagraph:P477:end

\end{minipage}\par\addvspace{4pt}

\reserve{7\baselineskip}
% mdparagraph:P478:begin
\section*{13. Cumulative practice}
% mdparagraph:P478:end

\reserve{7\baselineskip}
% mdparagraph:P479:begin
\subsection*{Combining the constructions}
% mdparagraph:P479:end

% mdparagraph:P480:begin
The second task contains deliberate errors to diagnose and correct; its input is not a model of correct Old English.
% mdparagraph:P480:end

\reserve{5\baselineskip}
\par\begin{minipage}{\linewidth}\raggedright
% mdparagraph:P481:begin
\textbf{X01.}
% mdparagraph:P481:end

% item: NEW-U09-X01
% mdparagraph:P482:begin
Translate this connected account: The fisherman commanded the men to carry the fish (sing.). I saw them carry it. I know that they helped him. Use hātan and gesēon with plain infinitives, then an indicative factual content clause.
% mdparagraph:P482:end

\end{minipage}\par\addvspace{4pt}

\par\begin{minipage}{\linewidth}\raggedright
% mdparagraph:P483:begin
\textbf{X02.}
% mdparagraph:P483:end

% item: NEW-U09-X02
% mdparagraph:P484:begin
Repair the case relations in these intended sentences and explain each repair: (a) ‘I see him carrying the fish’, written Ic gesēo hē ðone fisc beran.; (b) ‘God made the man good’, written God geworhte ðone mann gōd. Keep the selected complements.
% mdparagraph:P484:end

\end{minipage}\par\addvspace{4pt}

\clearpage
% mdparagraph:P485:begin
\section*{Sources and editorial notes}
% mdparagraph:P485:end

\begingroup\fontsize{10}{12.8}\selectfont
\setlength{\parskip}{1.5pt}
% mdparagraph:P486:begin
The principal editions are Albert S. Cook, \emph{A First Book in Old English: Grammar, Reader, Notes, and Vocabulary}, second edition, revised and enlarged (Ginn, 1895), and \emph{Exercises in Old English} (Ginn, 1895). The complete II.9–12 prompts are on \emph{Exercises} pp. 4–6 /\allowbreak{} PDF pp. 12–14. Their wording and original grammar references are retained. The twelve examples, fisherman dialogue and lettered tasks are new teaching material.
% mdparagraph:P486:end

% mdparagraph:P487:begin
“The Choice of Occupations” is on \emph{First Book} p. 131 /\allowbreak{} PDF p. 147. The four selected turns occupy prose lines 7–15 counted from the top of the page, with printed marginal numbers 10 and 15 anchoring the count. Cook’s words, spelling and conjectural insertions are retained. Source line breaks and superscript note cues are replaced by paragraphs and the separately printed original notes, whose numbers and anchors identify their locations. The preceding goods catalogue and price exchange are outside this excerpt.
% mdparagraph:P487:end

% mdparagraph:P488:begin
The complement and wish constructions follow Cook §§164–165, 169, 172–173, 193–196 and 199–200, with the partitive and appositional contrast in §§151 and 154b. R. D. Fulk, \emph{An Introductory Grammar of Old English with an Anthology of Readings}, author-corrected open-access reprint (2020), §§44–45 and 58, corroborates the mood and infinitive distinctions; his note on þrēo, p. 47, supports neuter agreement with mixed-gender groups. The preference, trade and service frames have contextual witnesses on Cook pp. 130–133.
% mdparagraph:P488:end

% mdparagraph:P489:begin
The noun and adjective patterns use Cook pp. 27–43 and the individual vocabulary entries. The verbal systems use §§92–119 and 141, with the stem alternations in §§15–19 and 33–37. Complete tables combine these rules with the checked lexical stems; they do not claim that every cell occurs in the reading. Joseph and Elizabeth Mary Wright, \emph{Old English Grammar}, second edition (1914), §§521–530, corroborates the weak-I forms. The selected wīsra genitive plural applies Wright §259.3, p. 130, on simplifying doubled r after another consonant.
% mdparagraph:P489:end

% mdparagraph:P490:begin
Bosworth–Toller, \emph{An Anglo-Saxon Dictionary} (1898), corroborates the specific forms or constructions under āstīgan (p. 55b), bycgan (p. 137a–b), gestillan (p. 446a), sęcgan (p. 855b) and absolute slēan (p. 883a); the \emph{Supplement} (1921), hwā III (pp. 575a–b), supports dependent questions. The main dictionary’s hund entry (p. 566a), together with Cook’s numeral rules on pp. 46–48, supports the locally supplied 365-day expression. That expression and libbende gesceafta are constructed supports, not quotations from the reader.
% mdparagraph:P490:end

% mdparagraph:P491:begin
Compound hyphens in dictionary citations are omitted in the ordinary displayed headwords; running source spelling, including for-ðām and swā-swā, is retained. The comparative wīsra and regular inflected cells are rule-based forms. The edition’s sēceað remains in the source beside selected productive sēcað. PDF locators are one-based: the selected \emph{First Book} scan uses printed page +16, \emph{Exercises} +8, Wright +20 and Fulk +12; dictionary a/\allowbreak{}b locators identify printed columns.
% mdparagraph:P491:end

\endgroup
\end{document}
